\documentclass[11pt, a4paper]{report}
% !TeX root = Lie Theory.tex
%----------------------------------------------------------------------------------------
%	PACKAGES (宏包管理)
%----------------------------------------------------------------------------------------
% 1. 编码与语言（中文支持）
\usepackage[UTF8, scheme=plain]{ctex}
% scheme=plain 保持定理/章节等名称为英文；若需全中文改为 scheme=chinese
% 2. 页面布局与排版
\usepackage{geometry}
% 不设 textheight，由 a4 页高与上下边距自动算出，避免 v 方向过指定
\geometry{
    textwidth=6in, % 稍微加宽以便放入更好的box
    top=1in,
    bottom=1in,
    headheight=12pt,
    headsep=25pt,
    footskip=30pt
}
\usepackage{setspace}
\setstretch{1.2}
\usepackage{float}
\usepackage{enumitem} % 更好的列表控制
\usepackage{tocloft}  % 加载目录控制宏包
\renewcommand{\cftsecleader}{\cftdotfill{\cftdotsep}} % 强制让 section 后面显示点号
% 3. 数学支持
\usepackage{amsmath, amsthm, amssymb, amsfonts}
\usepackage{mathtools}
\usepackage{thmtools} % 定理环境增强
\usepackage{extarrows}
% 4. 图形与颜色
\usepackage{graphicx}
\usepackage[dvipsnames, svgnames]{xcolor} % 加载更多颜色名
\usepackage{tikz}
\usepackage{tikz-cd}
\usepackage{ytableau} % 杨表
\usetikzlibrary{decorations.pathreplacing, calc, shadows}
% 5. 漂亮的文本框
\usepackage{tcolorbox}
\tcbuselibrary{breakable, skins, theorems}
% 6. 参考文献
%\usepackage[backend=biber,style=gb7714-2015,sorting=gb7714-2015]{biblatex}
%\addbibresource{references.bib}
% 7. 交互 (最后加载)
\usepackage[hidelinks]{hyperref}
\usepackage{array}
%----------------------------------------------------------------------------------------
%	CUSTOM COLORS (自定义配色方案)
%----------------------------------------------------------------------------------------
% 定理色系 - 蓝色
\definecolor{ThmColor}{RGB}{0, 71, 171}        % 深钴蓝
\definecolor{ThmBg}{RGB}{240, 244, 255}        % 极浅蓝
% 定义色系 - 绿色
\definecolor{DefColor}{RGB}{0, 100, 0}         % 深绿
\definecolor{DefBg}{RGB}{240, 255, 240}        % 蜜瓜绿
% 引理/命题色系 - 橙色
\definecolor{LemColor}{RGB}{204, 85, 0}        % 焦橙色
\definecolor{LemBg}{RGB}{255, 248, 240}        % 极浅橙
% 示例/注记色系 - 灰色/紫色
\definecolor{ExColor}{RGB}{105, 105, 105}      % 暗灰
\definecolor{ExBg}{RGB}{250, 250, 250}         % 极浅灰
% 专题/话题色系 - 靛紫色
\definecolor{TopicColor}{RGB}{75, 0, 130}      % 靛蓝色 (Indigo)
\definecolor{TopicBg}{RGB}{248, 245, 255}      % 极浅紫
%----------------------------------------------------------------------------------------
%	THEOREM ENVIRONMENTS (环境美化核心)
%----------------------------------------------------------------------------------------
% 1. Theorem (蓝色风格)
\declaretheorem[name=Theorem, numberwithin=section]{theorem}
\tcolorboxenvironment{theorem}{
    enhanced jigsaw,
    colframe=ThmColor,       % 边框颜色
    colback=ThmBg,           % 背景颜色
    coltitle=white,          
    boxrule=0.5pt,           % 边框粗细
    leftrule=2pt,            % 左边框加粗
    sharp corners=downhill,  % 样式
    breakable,               % 允许跨页
    drop shadow=black!5!white % 轻微阴影
}
% Corollary 使用与 Theorem 相同的风格，但通常紧随其后
\declaretheorem[name=Corollary, numberlike=theorem]{corollary}
\tcolorboxenvironment{corollary}{
    enhanced jigsaw,
    colframe=ThmColor!80!white, 
    colback=ThmBg,
    boxrule=0pt, leftrule=2pt, arc=0pt, outer arc=0pt,
    breakable
}
% 2. Definition (绿色风格)
\declaretheorem[name=Definition, numberlike=theorem]{definition}
\tcolorboxenvironment{definition}{
    enhanced jigsaw,
    colframe=DefColor,
    colback=DefBg,
    boxrule=0.5pt, leftrule=2pt,
    sharp corners=downhill,
    breakable,
    drop shadow=black!5!white
}
% 3. Lemma & Proposition (橙色风格)
\declaretheorem[name=Lemma, numberlike=theorem]{lemma}
\declaretheorem[name=Proposition, numberlike=theorem]{proposition}
\tcolorboxenvironment{lemma}{
    enhanced jigsaw,
    colframe=LemColor,
    colback=LemBg,
    boxrule=0.5pt, leftrule=2pt,
    breakable
}
\tcolorboxenvironment{proposition}{
    enhanced jigsaw,
    colframe=LemColor,
    colback=LemBg,
    boxrule=0.5pt, leftrule=2pt,
    breakable
}
% 4. Example & Remark (简约风格)
\declaretheorem[name=Example, numberlike=theorem]{example}
\declaretheorem[name=Remark, numberlike=theorem]{remark}
\declaretheorem[name=Exercise, numberlike=theorem]{exercise}
\tcolorboxenvironment{example}{
    enhanced jigsaw,
    colframe=ExColor,
    colback=ExBg,
    leftrule=2pt, boxrule=0pt,
    frame hidden, % 隐藏其他三边
    breakable
}
% 5. Topic (靛紫色风格，带标题支持)
\declaretheorem[name=Topic, numberlike=theorem]{topic}
\tcolorboxenvironment{topic}{
    enhanced jigsaw,
    breakable,
    colframe=TopicColor,        % 边框为靛紫色
    colback=TopicBg,           % 背景为极浅紫
    boxrule=0.5pt,             % 细边框
    leftrule=2pt,              % 左侧加粗线
    sharp corners=downhill,    % 与你其他环境一致的切角样式
    drop shadow=black!5!white, % 轻微阴影
    before skip=15pt,
    after skip=15pt
}
%----------------------------------------------------------------------------------------
%	MATH MACROS (保留你的常用定义)
%----------------------------------------------------------------------------------------
\newcommand{\g}{\mathfrak{g}}  
\newcommand{\der}{\mathrm{d}} 
\newcommand{\h}{\mathfrak{h}} 
\newcommand{\n}{\mathfrak{n}}  
\newcommand{\U}{\mathcal{U}}  
\newcommand{\C}{\mathbb{C}} 
\newcommand{\R}{\mathbb{R}} 
\newcommand{\borel}{\mathfrak{b}}
\newcommand{\Ind}{\mathrm{Ind}} 
\newcommand{\Hom}{\mathrm{Hom}}  
\newcommand{\Ext}{\mathrm{Ext}}
\newcommand{\HRule}[1]{\rule{\linewidth}{#1}}
%----------------------------------------------------------------------------------------
%	TITLE PAGE
%----------------------------------------------------------------------------------------
\begin{document}
\begin{titlepage}
    \centering
    \vspace*{2cm}

    \textsc{\Large Course Notes}\\[0.5cm]

    \HRule{1.5pt} \\[0.5cm]
    { \huge \bfseries Lie Theory }\\[0.4cm]
    \HRule{1.5pt} \\[1.5cm]

    \vspace{3cm}

    % 人员信息区域：垂直排列并区分字号
    \begin{spacing}{1.8}
        {\Large \textbf{Instructor:}} \\
        {\Large Prof. Efilm \textsc{Zelmanov}} \\[0.8cm]

        {\normalsize \textit{Student:}} \\
        {\normalsize Leo}
    \end{spacing}

    \vfill
    {\large \today}
\end{titlepage}
\newpage
\tableofcontents
\newpage


%----------------------------------------------------------------------------------------
%    CONTENT STRUCTURE TEMPLATE
%----------------------------------------------------------------------------------------
\chapter{Lecture 1 Where Lie algebras come from}

\section{Introduction}

Lie algebras were originally introduced by Sophus Lie in the 1870s to study the concept of continuous transformation groups. Lie's motivation was to develop a theory for differential equations analogous to Galois theory for algebraic equations. He discovered that the infinitesimal transformations of a continuous group form a linear space equipped with a bracket operation, now known as a Lie algebra. This linearization allows complex global geometric problems to be translated into algebraic problems, making them more tractable.

Let $ F $ be a field (mostly $ \operatorname{char}F =0$, even $ F=\mathbb{C} $).

\begin{definition}
    A linear algebra $ A $ is a vector space over $ F $, with a binary bilinear operation:
    \begin{itemize}
        \item binary operation: $ A \times A \to A $, $ (a,b) \mapsto ab $.
        \item bilinear: this operation is linear in each variable. In other word, $ \forall \alpha \in F $ and $ a,a',a'',b,b',b'' \in A $, we have:
              \begin{align*}
                  (\alpha a)b=a(\alpha b) = \alpha (ab) \\
                  (a'+a'')b = a'b+a''b                  \\
                  a(b'+b'')=ab'+ab''
              \end{align*}
    \end{itemize}
    The algebra $ A $ is associative if $ \forall a,b,c \in A $, $ (ab)c=a(bc) $.
\end{definition}

\begin{example}
    $ M_n(F) $, the algebra of $ n \times n $ matrices over $ F $.
\end{example}

\begin{example}
    Let $ V $ be a vector space, $ \operatorname{End}_F (V) =\{ \text{linear transformations } \varphi: V\to V \}$ is an associative algebra with the composition operation$ (\phi \psi) (v)=\phi(\psi(v))$.

    If $ \dim V < \infty $, when we fix a basis, then we fet a matrix representation of $ \varphi $ in our fix basis.
\end{example}

The most important thing everyone should know is Wedderburn's theorem.

\begin{theorem}[Wedderburn-Artin]
    Let $A$ be a semisimple Artinian ring (with 1). Then $A$ is a direct sum of matrix rings over division rings.

    Specifically, if ${}_A A \simeq S_1^{n_1} \oplus \dots \oplus S_r^{n_r}$ where $S_1, \dots, S_r$ are non-isomorphic simple modules occurring with multiplicities $n_1, \dots, n_r$, then
    \[
        A \simeq M_{n_1}(D_1) \oplus \dots \oplus M_{n_r}(D_r) \quad
    \]
    where $D_i = \operatorname{End}_A(S_i)^{op}$. This is a ring isomorphism, and each $M_{n_i}(D_i)$ is a 2-sided ideal.
\end{theorem}

\begin{proof}
    First, since $A$ is a semisimple ring, its left regular module ${}_A A$ can be decomposed into a direct sum of simple left $A$-modules:
    \[
        {}_A A \simeq S_1^{n_1} \oplus \dots \oplus S_r^{n_r} \quad
    \]
    where $S_1, \dots, S_r$ are pairwise non-isomorphic simple modules.

    By Schur's lemma, the endomorphism ring of a simple module is a division ring, so $D_i = \operatorname{End}_A(S_i)^{op}$ is a division ring.

    Next, we compute the endomorphism ring of the left regular module $\operatorname{End}_A({}_A A)$. On one hand, for any ring $A$ with 1, we have the ring isomorphism $\operatorname{End}_A({}_A A) \simeq A^{op}$. On the other hand, using our direct sum decomposition:
    \[
        \operatorname{End}_A({}_A A) \simeq \operatorname{End}_A(S_1^{n_1}) \oplus \dots \oplus \operatorname{End}_A(S_r^{n_r}) \quad
    \]
    This decomposition holds because $\operatorname{Hom}_A(S_i^{n_i}, S_j^{n_j}) = 0$ whenever $i \neq j$.

    For each component in the direct sum, we can rewrite the endomorphism ring using matrices over $D_i$:
    \[
        \operatorname{End}_A(S_i^{n_i}) \simeq M_{n_i}(D_i)^{op} \simeq M_{n_i}(D_i^{op}) \quad
    \]
    Substituting this back into our equation for $A^{op}$, we obtain:
    \[
        A^{op} \simeq M_{n_1}(D_1^{op}) \oplus \dots \oplus M_{n_r}(D_r^{op})
    \]

    To recover the ring $A$, we take the opposite ring on both sides. Recall the property that $M_n(R)^{op} \simeq M_n(R^{op})$. Applying this to each block yields:
    \[
        A \simeq (A^{op})^{op} \simeq \left( \bigoplus_{i=1}^r M_{n_i}(D_i^{op}) \right)^{op} \simeq \bigoplus_{i=1}^r M_{n_i}(D_i)
    \]
    Thus, the isomorphism is established.
\end{proof}

\begin{remark}
    If the base field $F$ is algebraically closed, every finite-dimensional
    division algebra over $F$ is equal to $F$. Hence the theorem simplifies to
    \[
        A \cong \prod_{i=1}^{r} M_{n_i}(F).
    \]
\end{remark}

\begin{example}
    Let $G$ be a finite group and let $\mathbb{C}[G]$ be its group algebra.
    By Maschke's theorem $\mathbb{C}[G]$ is semisimple. Therefore
    \[
        \mathbb{C}[G] \cong \prod_{i=1}^{r} M_{n_i}(\mathbb{C}),
    \]
    where $n_i$ are the dimensions of the irreducible complex representations
    of $G$.
\end{example}

\begin{definition}
    Suppose $ F $ is algebraically closed. An algebra $ A $ is simple if
    \begin{enumerate}
        \item $ \nexists \  I \triangleleft A $ and $ I \neq 0 $.
        \item $ A\cdot  A \neq (0)$.
    \end{enumerate}
\end{definition}

\begin{corollary}
    Any simple finite dimensional associative algebra is isomorphic to $ M_n(F) $ for some $ n $.
\end{corollary}

\begin{definition}[Jacobson Radical]
    Let $A$ be an associative algebra with identity over a field $F$.
    The \emph{Jacobson radical} of $A$, denoted by $\operatorname{Rad}(A)$
    or $J(A)$, is defined as the intersection of all maximal left ideals of $A$:
    \[
        \operatorname{Rad}(A)
        =
        \bigcap_{M \text{ maximal left ideal of } A} M.
    \]
\end{definition}


\begin{proposition}
    Let $A$ be an associative algebra with identity. Then
    \[
        \operatorname{Rad}(A)
        =
        \bigcap_{S \text{ simple left } A\text{-module}}
        \operatorname{Ann}_A(S),
    \]
    where
    \[
        \operatorname{Ann}_A(S)
        =
        \{a\in A \mid aS=0\}
    \]
    is the annihilator of $S$ in $A$.
\end{proposition}


\begin{proposition}[Equivalent Characterization]
    Let $A$ be an associative algebra with identity.
    Then an element $a\in A$ lies in $\operatorname{Rad}(A)$
    if and only if for every $x\in A$ the element
    \[
        1-ax
    \]
    is invertible in $A$.
\end{proposition}


\begin{proposition}
    The Jacobson radical $\operatorname{Rad}(A)$ is a two-sided ideal of $A$.
\end{proposition}


\begin{definition}[Semisimple Algebra]
    An associative algebra $A$ is called \emph{semisimple} if
    \[
        \operatorname{Rad}(A)=0.
    \]
\end{definition}

\section{Where do Lie algebras come from?}

\begin{definition}
    A Lie algebra is a vector space $ L $ over a field $ F $, with a binary operation $ [a,b] \in L $, such that
    \begin{enumerate}
        \item $ [a,b] $ is bilinear.
        \item $ [a,b] = -[b,a] $.
        \item $ [a,[b,c]] + [c,[a,b]] + [b,[c,a]] = 0 $, this is called Jacobi identity.
    \end{enumerate}
\end{definition}

\subsection{They come from associative algebras}

Let $ A $ be an associative algebra. $ \forall a,b \in A $, define $ [a,b]=ab-ba $. Check that $ (A,[ \ ,\  ]) $ is a Lie algebra:

\begin{proof}
    Bilinearity and skew-symmetry are clear. For Jacobi identity:
    \begin{align*}
        [a,[b,c]] & + [b,[c,a]] + [c,[a,b]]                                                 \\
                  & = a(bc-cb) - (bc-cb)a + b(ca-ac) - (ca-ac)b + c(ab-ba) - (ab-ba)c       \\
                  & = abc - acb - bca + cba + bca - bac - cab + acb + cab - cba - abc + bac \\
                  & = 0.
    \end{align*}
\end{proof}

We denote this new algebra $ (A,[ \ ,\  ]) $ by $ A^{(-)} $.

Suppose that $ L \subseteq A $ and $ [L,L] \subseteq L $. Then $ L \le A^{(-)} $。

\begin{example}
    $ M_n(F)^{(-)}=\mathfrak{gl}(n) $. We call it the general linear Lie algebra.

    $ [L,L]=\{ \sum_{i} [a_i,b_i] \} \trianglelefteq L $. Clearly $ [\mathfrak{gl}(n),\mathfrak{gl}(n)] \neq \mathfrak{gl}(n) $ and for $ n \ge 2 $, $ [\mathfrak{gl}(n),\mathfrak{gl}(n)]\neq 0 $.
\end{example}

\begin{example}
    Focus on $ [\mathfrak{gl}(n),\mathfrak{gl}(n)] $, all matrices are of trace 0. This is called the special linear Lie algebra, denoted by $ \mathfrak{sl}(n) $.
\end{example}

\begin{example}
    If the base field changed, some properties of the Lie algebra may change. For example, if $\operatorname{char}F=2$, then $[a,b]=[b,a]$. The Lie bracket is no longer skew-symmetric but actually symmetric.

    Consider $ [\mathfrak{gl}(n),\mathfrak{gl}(n)]=\mathfrak{sl}(n) $. If $\operatorname{char}F=0$ or $\operatorname{char}F \nmid n$, then $\mathfrak{sl}(n)$ is a simple Lie algebra. If $\operatorname{char}F\mid n$, then $\mathfrak{sl}(n)$ is not simple.
\end{example}

\begin{definition}
    $ Z $ is the center of Lie algebra $ L $ if $ Z=\{ z \in L | [z,L]=0 \} $.
\end{definition}

\subsection{They come from involutions of central simple algebras}

We can also find Lie algebras from the involution of a central simple algebra. Let $ A $ be central simple algebra and $ *: A \to A $ be a linear transformation. $ * $ is called a an involution if $ (a^*)^* =a $ and $ (ab)^*=b^*a^* $, $ \forall a,b \in A $.

\begin{proposition}
    Let $A$ be a central simple algebra over a field $F$ equipped with an involution
    $* : A \to A$. Then the set
    \[
        L=\{x\in A \mid *(x)=-x\}
    \]
    forms a Lie algebra under the commutator bracket
    \[
        [x,y]=xy-yx .
    \]
\end{proposition}

\begin{proof}
    First observe that any associative algebra $A$ becomes a Lie algebra
    with the commutator bracket
    \[
        [x,y]=xy-yx .
    \]

    Let
    \[
        L=\{x\in A\mid *(x)=-x\}.
    \]
    We show that $L$ is closed under the Lie bracket.

    Take $x,y\in L$. Then
    \[
        *(x)=-x,\qquad *(y)=-y.
    \]

    Since $*$ is an involution, it satisfies
    \[
        *(xy)=*(y)*(x).
    \]

    Hence
    \[
        *([x,y])
        =*(xy-yx)
        =*(y)*(x)-*(x)*(y).
    \]

    Substituting $*(x)=-x$ and $*(y)=-y$, we obtain
    \[
        *([x,y])
        =(-y)(-x)-(-x)(-y)
        =yx-xy
        =-[x,y].
    \]

    Thus $[x,y]\in L$. Therefore $L$ is closed under the commutator bracket.

    Since the commutator bracket in any associative algebra satisfies
    bilinearity, antisymmetry, and the Jacobi identity,
    it follows that $L$ is a Lie algebra.
\end{proof}

This construction explains why many classical Lie algebras arise from
central simple algebras with involution.  In particular, the classical
Lie algebras of Cartan--Killing type can be obtained in this way.

\begin{example}[Orthogonal involution]
    Let
    \[
        A=M_n(F)
    \]
    and define the involution
    \[
        *(X)=X^{T}.
    \]
    Then
    \[
        L=\{X\in M_n(F)\mid X^T=-X\}.
    \]
    This is the Lie algebra of skew-symmetric matrices, which is the
    orthogonal Lie algebra
    \[
        \mathfrak{so}_n(F).
    \]
\end{example}

\begin{example}[Symplectic involution]
    Let
    \[
        A=M_{2n}(F)
    \]
    and let
    \[
        J=
        \begin{pmatrix}
            0    & I_n \\
            -I_n & 0
        \end{pmatrix}.
    \]
    Define
    \[
        *(X)=J^{-1}X^{T}J.
    \]
    Then
    \[
        L=\{X\in M_{2n}(F)\mid *(X)=-X\}
    \]
    is the symplectic Lie algebra
    \[
        \mathfrak{sp}_{2n}(F).
    \]
\end{example}

\begin{example}[Unitary involution]
    Let $K/F$ be a quadratic field extension with Galois involution
    $a\mapsto \overline{a}$.  Consider
    \[
        A=M_n(K)
    \]
    with involution
    \[
        *(X)=\overline{X}^{\,T}.
    \]
    Then
    \[
        L=\{X\in M_n(K)\mid *(X)=-X\}
    \]
    is the unitary Lie algebra
    \[
        \mathfrak{su}_n .
    \]
\end{example}

Therefore classical Lie algebras arise from pairs $(A,*)$ where
$A$ is a central simple algebra and $*$ is an involution.  This
provides an algebraic framework for the Cartan--Killing classification
of classical Lie algebras.

\begin{example}
    The Quaternion algebra is an associative algebra with the involution $ a \mapsto a^* $. It is 4 dimensional central simple algebra with standard basis $ 1,i,j,k $, where the relations are $ i^2 = j^2 = k^2  = -1, ij=k, jk=i, ki=j $ as the following graph:
    \[
        \begin{tikzcd}
            i \arrow[r] & j \arrow[d]  \\
            & k \arrow[lu]
        \end{tikzcd}
    \]
    We use $ \mathbb{H} $ to denote the Quaternion algebra when we choose $ \mathbb{R} $ as the base field. i.e. $ \mathbb{H}=\{ a_0 1 + a_1 i + a_2 j + a_3 k | a_0,a_1,a_2,a_3 \in \mathbb{R} \} $. $ \forall x \in \mathbb{H} $ where $ x=a_0 1 + a_1 i + a_2 j + a_3 k $, $ * $ is defined as $ x^* = a_0 1 - a_1 i - a_2 j - a_3 k $ and $ x^{-1}= \frac{1}{a_0^2 + a_1^2 + a_2^2 + a_3^2} (a_0 1 - a_1 i - a_2 j - a_3 k) $. Hence $ \mathbb{H} $ is a division algebra.

    If the base field is algebraically closed, $F =\overline{F}$. We can denote the Quaternion algebra by $ 2 \times 2 $ matrices
    \[
        \begin{pmatrix}
            \alpha & \beta \\
            \gamma & \xi   \\
        \end{pmatrix}
    \]
    Where $ \alpha, \beta, \gamma, \xi \in F $. In this case it is no longer a division algebra ($a_0^2 + a_1^2 + a_2^2 + a_3^2$ may equal to 0), but just $ M_2(F) $. And in this case the involution $ * $ is defined as:
    \[
        \begin{pmatrix}
            \alpha & \beta \\
            \gamma & \xi   \\
        \end{pmatrix}
        \mapsto
        \begin{pmatrix}
            \xi     & -\beta \\
            -\gamma & \alpha \\
        \end{pmatrix}
    \]
    Which is the Symplectic involution.
\end{example}

\begin{example}
    From what we discussed. For a $2n \times 2n$ matrix $ A =\begin{pmatrix} A_1 & A_2 \\ A_3 & A_4 \end{pmatrix}$ in which $ A_1, A_2, A_3, A_4 $ are $ n \times n $ matrices, the unitary involution is defined as:
    \[
        A^*=
        \begin{pmatrix}
            A_4^t  & -A_3^t \\
            -A_2^t & A_1^t  \\
        \end{pmatrix}
    \]
    And the unitary involution is defined as:
    \[
        A^*=\overline{A}^t
    \]
    When $n=1$ we have two different type, and we will get two different Lie algebras. But they are isomorphic.
\end{example}

\begin{proposition}
    Let $F$ be a field of characteristic $\neq 2$, and let
    \[
        J=\begin{pmatrix}0&1\\-1&0\end{pmatrix}\in M_{2}(F).
    \]
    Define an involution $*$ on $M_{2}(F)$ by
    \[
        X^{*}=J^{-1}X^{t}J .
    \]
    Then the Lie algebra of $*$-skew elements,
    \[
        K(M_{2}(F),*)=\{X\in M_{2}(F)\mid X^{*}=-X\},
    \]
    is exactly $\mathfrak{sl}_{2}(F)$.
\end{proposition}

\begin{proof}
    Let
    \[
        X=\begin{pmatrix}a&b\\ c&d\end{pmatrix}\in M_{2}(F).
    \]
    Since
    \[
        J^{-1}=-J=\begin{pmatrix}0&-1\\1&0\end{pmatrix},
    \]
    we compute
    \[
        X^{*}
        =J^{-1}X^{t}J
        =\begin{pmatrix}0&-1\\1&0\end{pmatrix}
        \begin{pmatrix}a&c\\ b&d\end{pmatrix}
        \begin{pmatrix}0&1\\-1&0\end{pmatrix}
        =\begin{pmatrix}d&-b\\-c&a\end{pmatrix}.
    \]
    Now the condition $X^{*}=-X$ becomes
    \[
        \begin{pmatrix}d&-b\\-c&a\end{pmatrix}
        =
        \begin{pmatrix}-a&-b\\-c&-d\end{pmatrix}.
    \]
    Hence $d=-a$, while $b$ and $c$ are arbitrary. Therefore
    \[
        K(M_{2}(F),*)
        =
        \left\{
        \begin{pmatrix}
            a & b  \\
            c & -a
        \end{pmatrix}
        \,\middle|\,
        a,b,c\in F
        \right\}.
    \]
    But this is precisely the set of all $2\times 2$ matrices of trace $0$, namely
    \[
        \mathfrak{sl}_{2}(F)
        =
        \left\{
        \begin{pmatrix}
            a & b  \\
            c & -a
        \end{pmatrix}
        \,\middle|\,
        a,b,c\in F
        \right\}.
    \]
    Thus
    \[
        K(M_{2}(F),*)=\mathfrak{sl}_{2}(F).
    \]
    Since $K(M_{2}(F),*)$ is closed under the commutator bracket, it follows that the Lie algebra obtained from the symplectic involution on $M_{2}(F)$ is $\mathfrak{sl}_{2}(F)$.
\end{proof}

\begin{example}
    Let $K/F$ be a quadratic field extension, and let $a \mapsto \overline{a}$ denote the nontrivial
    $F$-automorphism of $K$. Consider the algebra $M_2(K)$ and define an involution $*$ by
    \[
        X^* = \overline{X}^{\,t}.
    \]
    Then $*$ is a unitary involution on $M_2(K)$.

    The Lie algebra of $*$-skew elements is
    \[
        K(M_2(K),*)=\{X\in M_2(K)\mid X^*=-X\}.
    \]
    If
    \[
        X=\begin{pmatrix} a & b \\ c & d \end{pmatrix},
    \]
    then
    \[
        X^*=\overline{X}^{\,t}
        =
        \begin{pmatrix}
            \overline{a} & \overline{c} \\
            \overline{b} & \overline{d}
        \end{pmatrix}.
    \]
    Thus the condition $X^*=-X$ is equivalent to
    \[
        \overline{a}=-a,\qquad \overline{d}=-d,\qquad \overline{c}=-b.
    \]
    Hence every element of $K(M_2(K),*)$ has the form
    \[
        X=
        \begin{pmatrix}
            a             & b \\
            -\overline{b} & d
        \end{pmatrix},
        \qquad \text{with } \overline{a}=-a,\ \overline{d}=-d.
    \]
    This is the unitary Lie algebra associated with the standard hermitian form.

    Its derived Lie algebra consists of those elements with trace $0$, namely
    \[
        \mathfrak{su}_2(K/F)
        =
        \left\{
        \begin{pmatrix}
            a             & b  \\
            -\overline{b} & -a
        \end{pmatrix}
        \,\middle|\,
        a,b\in K,\ \overline{a}=-a
        \right\}.
    \]
    After extending scalars to a splitting field (for instance to an algebraic closure of $F$),
    this Lie algebra becomes isomorphic to $\mathfrak{sl}_2$. Therefore the Lie algebra arising
    from a unitary involution in the $2\times 2$ case is of type $A_1$.
\end{example}

\begin{remark}
    Why does this Lie algebra become $\mathfrak{sl}_2$? The reason is that unitary Lie algebras are
    forms of special linear Lie algebras. In the present $2\times 2$ case, the corresponding type is
    $A_1$. Over a field where the quadratic extension splits, the involution becomes equivalent to the
    ordinary transpose situation, and the resulting special unitary Lie algebra becomes the split simple
    Lie algebra of type $A_1$, namely
    \[
        \mathfrak{sl}_2.
    \]
    Thus, although over the ground field one naturally writes the Lie algebra as $\mathfrak{su}_2$,
    its split form is precisely $\mathfrak{sl}_2$.
\end{remark}

\begin{example}[The unitary involution on $M_{2}(\mathbb{C})$ and the Lie algebra $\mathfrak{su}_{2}$]
    Let $F=\mathbb{R}$ and $K=\mathbb{C}$, and let $\overline{\phantom{x}}$ denote complex conjugation.
    Consider the algebra $M_{2}(\mathbb{C})$ equipped with the involution
    \[
        X^{*}=\overline{X}^{\,t}.
    \]
    This is the standard unitary involution on $M_{2}(\mathbb{C})$ relative to the extension
    $\mathbb{C}/\mathbb{R}$.

    We first determine the Lie algebra of $*$-skew elements:
    \[
        K(M_{2}(\mathbb{C}),*)=\{X\in M_{2}(\mathbb{C})\mid X^{*}=-X\}.
    \]
    Let
    \[
        X=\begin{pmatrix}a&b\\ c&d\end{pmatrix}\in M_{2}(\mathbb{C}).
    \]
    Then
    \[
        X^{*}
        =
        \overline{X}^{\,t}
        =
        \begin{pmatrix}
            \overline{a} & \overline{c} \\
            \overline{b} & \overline{d}
        \end{pmatrix}.
    \]
    Thus the condition $X^{*}=-X$ is equivalent to
    \[
        \overline{a}=-a,\qquad \overline{d}=-d,\qquad \overline{c}=-b.
    \]
    Hence $a$ and $d$ are purely imaginary, and $c=-\overline{b}$. Therefore
    \[
        K(M_{2}(\mathbb{C}),*)
        =
        \left\{
        \begin{pmatrix}
            a             & b \\
            -\overline{b} & d
        \end{pmatrix}
        \,\middle|\,
        a,d\in i\mathbb{R},\ b\in\mathbb{C}
        \right\}.
    \]
    This is the unitary Lie algebra $\mathfrak{u}_{2}$.

    Its derived Lie algebra is the trace-zero part, namely
    \[
        \mathfrak{su}_{2}
        =
        \left\{
        \begin{pmatrix}
            a             & b  \\
            -\overline{b} & -a
        \end{pmatrix}
        \,\middle|\,
        a\in i\mathbb{R},\ b\in\mathbb{C}
        \right\}.
    \]
    Indeed, $\operatorname{tr}([X,Y])=0$ for all $X,Y\in M_{2}(\mathbb{C})$, so every commutator in
    $\mathfrak{u}_{2}$ has trace $0$. Conversely, in this $2\times2$ case, the trace-zero subspace is
    generated by commutators, and hence it is exactly the derived Lie algebra.

    To see that $\mathfrak{su}_{2}$ is $3$-dimensional over $\mathbb{R}$, write
    \[
        a=it,\qquad b=u+iv,
        \qquad t,u,v\in\mathbb{R}.
    \]
    Then every element of $\mathfrak{su}_{2}$ has the form
    \[
        \begin{pmatrix}
            it    & u+iv \\
            -u+iv & -it
        \end{pmatrix},
        \qquad t,u,v\in\mathbb{R}.
    \]
    Therefore
    \[
        \dim_{\mathbb{R}}\mathfrak{su}_{2}=3.
    \]
    A convenient real basis is
    \[
        H=
        \begin{pmatrix}
            i & 0  \\
            0 & -i
        \end{pmatrix},
        \qquad
        X=
        \begin{pmatrix}
            0  & 1 \\
            -1 & 0
        \end{pmatrix},
        \qquad
        Y=
        \begin{pmatrix}
            0 & i \\
            i & 0
        \end{pmatrix}.
    \]
    One checks directly that
    \[
        [H,X]=2Y,\qquad [H,Y]=-2X,\qquad [X,Y]=2H,
    \]
    so $\mathfrak{su}_{2}$ is a $3$-dimensional real Lie algebra.

    Now we explain why this Lie algebra becomes $\mathfrak{sl}_{2}(\mathbb{C})$ after extending scalars
    from $\mathbb{R}$ to $\mathbb{C}$. Consider the complexification
    \[
        \mathfrak{su}_{2}\otimes_{\mathbb{R}}\mathbb{C}.
    \]
    Since $H,X,Y$ form a real basis of $\mathfrak{su}_{2}$, they also form a complex basis of
    $\mathfrak{su}_{2}\otimes_{\mathbb{R}}\mathbb{C}$. Define
    \[
        h=-iH=
        \begin{pmatrix}
            1 & 0  \\
            0 & -1
        \end{pmatrix},
        \
        e=\frac{1}{2}(X-iY)=
        \begin{pmatrix}
            0 & 1 \\
            0 & 0
        \end{pmatrix},
        \
        f=\frac{1}{2}(-X-iY)=
        \begin{pmatrix}
            0  & 0 \\
            -1 & 0
        \end{pmatrix}.
    \]
    These elements lie in $\mathfrak{su}_{2}\otimes_{\mathbb{R}}\mathbb{C}$ and satisfy
    \[
        [h,e]=2e,\qquad [h,f]=-2f,\qquad [e,f]=h.
    \]
    Thus they form a standard $\mathfrak{sl}_{2}$-triple, and therefore
    \[
        \mathfrak{su}_{2}\otimes_{\mathbb{R}}\mathbb{C}\cong \mathfrak{sl}_{2}(\mathbb{C}).
    \]

    In other words, $\mathfrak{su}_{2}$ is a real form of $\mathfrak{sl}_{2}(\mathbb{C})$. Hence the Lie
    algebra arising from the unitary involution on $M_{2}(\mathbb{C})$ is of type $A_{1}$.
\end{example}

\begin{remark}
    This explains why, in the $2\times 2$ case, the Lie algebra arising from a symplectic involution is isomorphic to the one of type $A_{1}$. Equivalently, in this lowest-rank case one has the accidental isomorphism
    \[
        C_{1}\cong A_{1}.
    \]
\end{remark}

\subsection{They come from the derivation of any algebra}

Let $A$ be any algebra.

\begin{definition}
    A derivation of $A$ is a linear map $\mathrm{d}: A \to A$ such that $\mathrm{d}(ab)=\mathrm{d}(a)b+a\mathrm{d}(b)$ for all $a,b \in A$.
\end{definition}

\begin{example}
    Let $A=F[t_1,\ldots ,t_n]$ be the polynomial algebra. We have the formal derivative $ \partial $ where $\partial(t_i^n)=nt_i^{n-1}$ and $\partial(fg)=f\partial(g)+\partial(f)g$.

    Let $\mathrm{d}=\sum_{i=1}^n f_i(t_1,\ldots ,t_n) \frac{\partial}{\partial t_i} $, this is a derivation of $A$.
    \[
        \mathrm{d}(f)=\sum_{i=1}^n f_i(t_1,\ldots ,t_n) \frac{\partial f}{\partial t_i}
    \]
    If $\mathrm{d}_1,\mathrm{d}_2$ are derivations of $A$, then $[\mathrm{d}_1,\mathrm{d}_2]$ is again a derivation of $A$. We can check it by the following calculation:
    \begin{align*}
        [\mathrm{d}_1,\mathrm{d}_2](ab) & =\mathrm{d}_1(\mathrm{d}_2(ab))-\mathrm{d}_2(\mathrm{d}_1(ab))                                                               \\
                                        & =\mathrm{d}_1(\mathrm{d}_2(a)b+a\mathrm{d}_2(b))-\mathrm{d}_2(\mathrm{d}_1(a)b+a\mathrm{d}_1(b))                             \\
                                        & =\mathrm{d}_1(\mathrm{d}_2(a))b+\mathrm{d}_1(a\mathrm{d}_2(b))+\mathrm{d}_2(\mathrm{d}_1(a))b+\mathrm{d}_2(a\mathrm{d}_1(b)) \\
                                        & =\mathrm{d}_1(\mathrm{d}_2(a))b+\mathrm{d}_2(\mathrm{d}_1(a))b+\mathrm{d}_1(a\mathrm{d}_2(b))+\mathrm{d}_2(a\mathrm{d}_1(b)) \\
                                        & =\mathrm{d}_1(\mathrm{d}_2(a))b+\mathrm{d}_2(\mathrm{d}_1(a))b+a[\mathrm{d}_1,\mathrm{d}_2](b)                               \\
                                        & =[\mathrm{d}_1,\mathrm{d}_2](a)b+a[\mathrm{d}_1,\mathrm{d}_2](b)                                                             \\
                                        & =[\mathrm{d}_1,\mathrm{d}_2](ab)
    \end{align*}
    Therefore $[\mathrm{d}_1,\mathrm{d}_2]$ is a derivation of $A$.
\end{example}

The above example means that:
\[
    \operatorname{Der}(A)=\{\ \mathrm{d}: A \to A \mid \mathrm{d} \text{ is a derivation} \} < \operatorname{End}_F(A)^{(-)}
\]
In particular, when $A=F[t_1,\ldots ,t_n]$, we have:
\[
    \operatorname{Der}(A)=\{\ \sum_{i=1}^n f_i(t_1,\ldots ,t_n) \frac{\partial}{\partial t_i} \mid f_i(t_1,\ldots ,t_n) \in A \} < \operatorname{End}_F(A)^{(-)}
\]
We denote this algebra by $W(n)$. It is called the $n$-th Witt algebra. It is a simple, infinite dimensional Lie algebra (when $\operatorname{char}F=0$).

Thus every element of $W(n)$ has the form
\[
    D=\sum_{i=1}^n f_i\,\partial_i,
    \qquad f_i\in A.
\]

\begin{definition}
    The divergence map is defined by
    \[
        \operatorname{div}:W(n)\to A,
        \qquad
        \operatorname{div}\!\left(\sum_{i=1}^n f_i\partial_i\right)
        =\sum_{i=1}^n \partial_i(f_i).
    \]
    Its kernel is the \emph{special} Lie algebra
    \[
        S(n):=\ker(\operatorname{div})
        =\left\{
        \sum_{i=1}^n f_i\partial_i\in W(n)
        \;\middle|\;
        \sum_{i=1}^n \partial_i(f_i)=0
        \right\}.
    \]
    Equivalently, $S(n)$ is the Lie algebra of polynomial vector fields preserving
    the standard volume form
    \[
        \omega=dx_1\wedge\cdots\wedge dx_n.
    \]
\end{definition}


A basic family of elements of $S(n)$ is given by
\[
    D_{ij}(f):=\partial_j(f)\partial_i-\partial_i(f)\partial_j,
    \qquad
    1\le i<j\le n,\quad f\in A,
\]
since
\[
    \operatorname{div}(D_{ij}(f))
    =
    \partial_i\partial_j(f)-\partial_j\partial_i(f)=0.
\]
One can check
\[
    S(n)=< D_{ij}(f) \mid 1\le i<j\le n, f\in A >_{Lie}
\]
These elements generate $S(n)$ as a Lie algebra.

\begin{definition}
    The grading of $W(n)$ is defined by
    Give $A$ its standard grading by declaring $\deg x_i=1$. Then $W(n)$ becomes a
    $\mathbb Z$-graded Lie algebra with
    \[
        \deg(x^a\partial_i)=|a|-1,
        \qquad
        |a|:=a_1+\cdots+a_n,
    \]
    so that
    \[
        W(n)=\bigoplus_{k\ge -1} W(n)_k.
    \]

    and
    \[
        W(n)_0
        =
        \operatorname{span}_F\{x_j\partial_i\mid 1\le i,j\le n\}
        \cong \mathfrak{gl}_n(F).
    \]
\end{definition}


Since the divergence map is homogeneous of degree $0$, the subalgebra $S(n)$ is
graded:
\[
    S(n)=\bigoplus_{k\ge -1} S(n)_k,
    \qquad
    S(n)_k=S(n)\cap W(n)_k.
\]
Moreover,
\[
    S(n)_{-1}=W(n)_{-1},
\]
and
\[
    S(n)_0
    =
    \left\{
    \sum_{i,j} a_{ij}x_j\partial_i
    \;\middle|\;
    \sum_i a_{ii}=0
    \right\}
    \cong \mathfrak{sl}_n(F),
\]
because
\[
    \operatorname{div}\!\left(\sum_{i,j} a_{ij}x_j\partial_i\right)
    =
    \sum_i a_{ii}.
\]

In characteristic $0$, the special algebra is perfect:
\[
    [S(n),S(n)]=S(n).
\]
Hence
\[
    S(n)^{(1)}=S(n).
\]
For $n\ge 2$, the Lie algebra $S(n)$ is simple. Therefore, in characteristic
$0$ one has
\[
    W(n)\supset S(n)=S(n)^{(1)},
\]
where both algebras are infinite-dimensional, $W(n)$ is the full Lie algebra of
polynomial vector fields, and $S(n)$ is its simple divergence-zero subalgebra.

\begin{remark}
    In characteristic $0$, the chain $W(n)\supset S(n)\supset S(n)^{(1)}$ is thus
    less informative than in positive characteristic, since the last inclusion is an
    equality.
\end{remark}

Let $A$ be an associative algebra. Fix an element $a \in A$, we define : $\mathrm{ad}(a): A \to A : b\mapsto [a,b]$. This is a derivation of $A$.

\begin{definition}
    Let $L$ be a Lie algebra. The adjoint map is defined as:
    \[
        \mathrm{ad}: L \to \operatorname{End}_F(L)^{(-)}, a \mapsto \mathrm{ad}(a)
    \]
    where $\mathrm{ad}(a): L \to L : b\mapsto [a,b]$ is an inner derivation of $L$. We also call $\mathrm{ad}(a)$ an adjoint operator.
\end{definition}

Note that $\mathrm{ad}(a) \in \operatorname{Der}(L)$ check it by the following calculation:
\begin{align*}
    \mathrm{ad}(a)([b,c]) & =\mathrm{ad}(a)(bc-cb)                       \\
                          & =[a,bc]-[a,cb]                               \\
                          & =[a,b]c+b[a,c]-[a,c]b-c[a,b]                 \\
                          & =[\mathrm{ad}(a)(b),c]+[b,\mathrm{ad}(a)(c)]
\end{align*}

"Inner" here comes from that $\mathrm{ad}(L) \triangleleft \operatorname{Der}(L)$ since $\forall \mathrm{d} \in \operatorname{Der}(L)$, we have $[\mathrm{d},\mathrm{ad}(a)]=  \mathrm{ad}(\mathrm{d}(a))$. Check it by the following calculation:
\begin{align*}
    [d,\operatorname{ad}(a)](x)
     & = d\bigl([a,x]\bigr)-[a,d(x)] \\
     & = [d(a),x]+[a,d(x)]-[a,d(x)]  \\
     & = [d(a),x]                    \\
     & = \operatorname{ad}(d(a))(x).
\end{align*}

\begin{remark}[Exponentials of nilpotent inner derivations]
    Let $L$ be a Lie algebra over a field $F$, and let
    \[
        d \in \operatorname{ad}(L) \subseteq \operatorname{Der}(L).
    \]
    Thus there exists $a \in L$ such that
    \[
        d=\operatorname{ad}(a), \qquad d(x)=[a,x] \quad \text{for all } x\in L.
    \]

    Assume that $d$ is nilpotent, so that $d^N=0$ for some integer $N\ge 1$.
    Then one defines the exponential of $d$ by
    \[
        \exp(d):=\sum_{n=0}^{\infty}\frac{d^n}{n!}.
    \]
    Since $d$ is nilpotent, this is actually a finite sum:
    \[
        \exp(d)=I+d+\frac{d^2}{2!}+\cdots+\frac{d^{N-1}}{(N-1)!}.
    \]
    Hence $\exp(d)$ is a well-defined linear endomorphism of $L$.

    Moreover, $\exp(d)$ is invertible, with inverse
    \[
        \exp(d)^{-1}=\exp(-d).
    \]
    Since $d$ is a derivation, $\exp(d)$ is in fact a Lie algebra automorphism:
    \[
        \exp(d)([x,y])=[\exp(d)(x),\exp(d)(y)]
        \qquad \text{for all } x,y\in L.
    \]
    Therefore
    \[
        \exp(d)\in \operatorname{Aut}(L).
    \]

    In particular, if $d=\operatorname{ad}(a)$, then
    \[
        \exp(\operatorname{ad}(a))(x)
        =
        x+[a,x]+\frac{1}{2!}[a,[a,x]]+\frac{1}{3!}[a,[a,[a,x]]]+\cdots.
    \]
    When $\operatorname{ad}(a)$ is nilpotent, this sum is finite. The automorphism
    \[
        \exp(\operatorname{ad}(a))
    \]
    is called an inner automorphism arising from the nilpotent inner derivation $\operatorname{ad}(a)$.

    A fundamental example is the matrix Lie algebra $L=\mathfrak{gl}_n(F)$. For $A\in \mathfrak{gl}_n(F)$, one has
    \[
        \operatorname{ad}(A)(X)=[A,X]=AX-XA,
    \]
    and
    \[
        \exp(\operatorname{ad}(A))(X)=e^A X e^{-A}.
    \]
    Thus, in the matrix case, the exponential of an inner derivation is exactly conjugation by the matrix exponential.
\end{remark}

\begin{proposition}
    Let $L$ be a Lie algebra over a field $F$, and let $d\in \operatorname{Der}(L)$ be nilpotent. Then
    \[
        \exp(d):=\sum_{n=0}^{\infty}\frac{d^n}{n!}
    \]
    is a Lie algebra automorphism of $L$.
\end{proposition}

\begin{proof}
    Since $d$ is nilpotent, there exists $N\ge 1$ such that $d^N=0$. Hence
    \[
        \exp(d)=I+d+\frac{d^2}{2!}+\cdots+\frac{d^{N-1}}{(N-1)!}
    \]
    is a finite sum, so it is a well-defined linear endomorphism of $L$.

    Moreover,
    \[
        \exp(d)\exp(-d)=\exp(-d)\exp(d)=I,
    \]
    because $d$ commutes with itself and the usual exponential identity holds for polynomials in $d$.
    Thus $\exp(d)$ is invertible, with inverse $\exp(-d)$.

    It remains to show that $\exp(d)$ preserves the Lie bracket. Define
    \[
        \phi_t:=\exp(td)=\sum_{n=0}^{\infty}\frac{t^n d^n}{n!},
    \]
    where the sum is again finite since $d$ is nilpotent. For $x,y\in L$, set
    \[
        F(t):=\phi_t([x,y]), \qquad G(t):=[\phi_t(x),\phi_t(y)].
    \]
    We show that $F(t)=G(t)$.

    First,
    \[
        F'(t)=d(\phi_t([x,y]))=\phi_t(d([x,y])).
    \]
    Since $d$ is a derivation,
    \[
        d([x,y])=[d(x),y]+[x,d(y)].
    \]
    Therefore
    \[
        F'(t)=\phi_t([d(x),y]+[x,d(y)])
        =[\phi_t(d(x)),\phi_t(y)]+[\phi_t(x),\phi_t(d(y))].
    \]
    On the other hand,
    \[
        G'(t)=[\phi_t'(x),\phi_t(y)]+[\phi_t(x),\phi_t'(y)]
        =[\phi_t(d(x)),\phi_t(y)]+[\phi_t(x),\phi_t(d(y))].
    \]
    Hence
    \[
        F'(t)=G'(t).
    \]
    Also,
    \[
        F(0)=[x,y]=G(0).
    \]
    It follows that $F(t)=G(t)$ for all $t$, and in particular at $t=1$ we get
    \[
        \exp(d)([x,y])=[\exp(d)(x),\exp(d)(y)].
    \]
    Thus $\exp(d)$ is a Lie algebra automorphism of $L$.
\end{proof}

\subsection{They come from brackets}

\begin{definition}[Poisson bracket]
    Let $A$ be a commutative associative algebra over a field $\Bbbk$.
    A \emph{Poisson bracket} on $A$ is a bilinear map
    \[
        \{\cdot,\cdot\}:A\times A\to A
    \]
    such that for all $f,g,h\in A$ the following conditions hold:
    \begin{enumerate}
        \item \textbf{Skew-symmetry:}
              \[
                  \{f,g\}=-\{g,f\}.
              \]
        \item \textbf{Jacobi identity:}
              \[
                  \{f,\{g,h\}\}+\{g,\{h,f\}\}+\{h,\{f,g\}\}=0.
              \]
        \item \textbf{Leibniz rule:}
              \[
                  \{f,gh\}=\{f,g\}h+g\{f,h\}.
              \]
    \end{enumerate}
    In other words, $(A,\{\cdot,\cdot\})$ is a Lie algebra, and for each fixed
    $f\in A$, the map
    \[
        \mathrm{ad}_f:A\to A,\qquad \mathrm{ad}_f(g)=\{f,g\},
    \]
    is a derivation of the algebra $A$.
    An algebra $A$ equipped with such a bracket is called a \emph{Poisson algebra}.
\end{definition}

\begin{remark}
    The Leibniz rule in the second variable follows automatically from skew-symmetry:
    \[
        \{fg,h\}=f\{g,h\}+g\{f,h\}.
    \]
    Thus a Poisson bracket is precisely a Lie bracket compatible with the
    commutative multiplication on $A$ by the derivation property.
\end{remark}

\begin{example}[Canonical Poisson bracket on $\mathbb{R}^{2n}$]
    Let
    \[
        A=C^\infty(\mathbb{R}^{2n}),
        \qquad
        (q_1,\dots,q_n,p_1,\dots,p_n)
    \]
    be the standard coordinates. Define
    \[
        \{f,g\}
        =
        \sum_{i=1}^n
        \left(
        \frac{\partial f}{\partial q_i}\frac{\partial g}{\partial p_i}
        -
        \frac{\partial f}{\partial p_i}\frac{\partial g}{\partial q_i}
        \right).
    \]
    Then $\{\cdot,\cdot\}$ is a Poisson bracket on $A$.
    In particular,
    \[
        \{q_i,q_j\}=0,\qquad
        \{p_i,p_j\}=0,\qquad
        \{q_i,p_j\}=\delta_{ij}.
    \]
\end{example}

\begin{definition}
    Let $F$ be a field of characteristic $0$, and let $\mathcal{A}$ be an algebra of differentiable functions in one variable $x$.
    For $f,g\in \mathcal{A}$, define
    \[
        [f,g]:=fg'-f'g,
    \]
    where $f'=\dfrac{df}{dx}$ and $g'=\dfrac{dg}{dx}$.
\end{definition}

\begin{definition}
    To each $f\in \mathcal{A}$, associate the vector field
    \[
        X_f:=f(x)\frac{d}{dx}.
    \]
\end{definition}

\begin{proposition}
    For $f,g\in \mathcal{A}$, one has
    \[
        [X_f,X_g]=X_{[f,g]}.
    \]
    Equivalently,
    \[
        [X_f,X_g]=\left(fg'-gf'\right)\frac{d}{dx}.
    \]
\end{proposition}

\begin{proof}
    For any test function $u$, we have
    \[
        X_f(u)=fu', \qquad X_g(u)=gu'.
    \]
    Thus
    \[
        X_f(X_g(u))=f(gu')'=fg'u'+fgu'',
    \]
    and similarly
    \[
        X_g(X_f(u))=g(fu')'=gf'u'+gfu''.
    \]
    Subtracting, we obtain
    \[
        [X_f,X_g](u)=X_f(X_g(u))-X_g(X_f(u))=(fg'-gf')u'.
    \]
    Therefore
    \[
        [X_f,X_g]=\left(fg'-gf'\right)\frac{d}{dx}=X_{[f,g]}.
    \]
\end{proof}

\begin{proposition}
    The operation
    \[
        [f,g]=fg'-f'g
    \]
    is bilinear and skew-symmetric.
\end{proposition}

\begin{proof}
    Bilinearity follows immediately from the linearity of differentiation. Also,
    \[
        [g,f]=gf'-g'f=-(fg'-f'g)=-[f,g].
    \]
    Hence the bracket is skew-symmetric.
\end{proof}

\begin{theorem}
    The bracket
    \[
        [f,g]=fg'-f'g
    \]
    defines a Lie algebra structure on $\mathcal{A}$.
\end{theorem}

\begin{proof}
    By the previous proposition, it remains only to verify the Jacobi identity. This follows from the fact that the map
    \[
        f\longmapsto X_f=f\frac{d}{dx}
    \]
    identifies this bracket with the usual Lie bracket of vector fields on the line, and the latter always satisfies the Jacobi identity. Hence
    \[
        [f,[g,h]]+[g,[h,f]]+[h,[f,g]]=0
    \]
    for all $f,g,h\in \mathcal{A}$.
\end{proof}

\subsection{They come from Lie groups}

\begin{definition}
    Let \(G\) be a Lie group with identity element \(e\). The \emph{Lie algebra} of \(G\), denoted by \(\operatorname{Lie}(G)\) or \(\mathfrak g\), is the tangent space
    \[
        \mathfrak g:=T_eG
    \]
    equipped with the bracket induced by left-invariant vector fields.
\end{definition}

\begin{definition}
    For each \(X\in T_eG\), the corresponding \emph{left-invariant vector field} \(\widetilde X\) on \(G\) is defined by
    \[
        \widetilde X_g:=(dL_g)_e(X),
        \qquad g\in G,
    \]
    where \(L_g:G\to G\) is left translation, \(L_g(h)=gh\).
\end{definition}

\begin{proposition}
    The assignment \(X\mapsto \widetilde X\) gives a bijection between \(T_eG\) and the space of left-invariant vector fields on \(G\).
\end{proposition}

\begin{definition}
    For \(X,Y\in \mathfrak g=T_eG\), define
    \[
        [X,Y]:=[\widetilde X,\widetilde Y](e),
    \]
    where on the right-hand side \([\widetilde X,\widetilde Y]\) is the usual Lie bracket of vector fields on \(G\).
\end{definition}

\begin{proposition}
    The bracket defined above makes \(T_eG\) into a Lie algebra. In particular, it is bilinear, alternating, and satisfies the Jacobi identity.
\end{proposition}

\begin{example}
    If \(G=GL_n(F)\), where \(F=\mathbb R\) or \(\mathbb C\), then
    \[
        \operatorname{Lie}(G)=T_IGL_n(F)\cong M_n(F),
    \]
    and the Lie bracket is
    \[
        [X,Y]=XY-YX.
    \]
    Hence
    \[
        \operatorname{Lie}(GL_n(F))=\mathfrak{gl}_n(F).
    \]
\end{example}

\begin{example}
    If \(G=SL_n(F)\), then
    \[
        \operatorname{Lie}(SL_n(F))=\mathfrak{sl}_n(F)
        =\{X\in M_n(F)\mid \operatorname{tr}(X)=0\}.
    \]
    Indeed, differentiating the defining equation \(\det(g)=1\) at the identity gives the trace-zero condition.
\end{example}

\begin{example}
    If \(G=O_n(\mathbb R)\), then
    \[
        \operatorname{Lie}(O_n(\mathbb R))=\mathfrak{so}_n(\mathbb R)
        =\{X\in M_n(\mathbb R)\mid X^T=-X\}.
    \]
    This follows by differentiating the relation \(g^Tg=I\) at the identity.
\end{example}

\section{Basic concepts of Lie algebra}

Let \(L\) be a Lie algebra over a field \(F\).

\begin{definition}
    The \emph{derived series} of \(L\) is defined inductively by
    \[
        L^{(0)}=L, \qquad L^{(1)}=[L,L], \qquad L^{(i)}=[L^{(i-1)},L^{(i-1)}]\quad (i\ge 1).
    \]
    The Lie algebra \(L\) is called \emph{solvable} if
    \[
        L^{(n)}=0
    \]
    for some integer \(n\ge 0\).
\end{definition}

\begin{definition}
    The \emph{descending central series} of \(L\) is defined by
    \[
        L^1=L, \qquad L^2=[L,L], \qquad L^i=[L,L^{i-1}]\quad (i\ge 2).
    \]
    The Lie algebra \(L\) is called \emph{nilpotent} if
    \[
        L^n=0
    \]
    for some integer \(n\ge 1\).
\end{definition}

\begin{proposition}
    If \(L\) is nilpotent, then \(L\) is solvable.
\end{proposition}

\begin{proof}
    It is enough to show that
    \[
        L^{(i)}\subseteq L^{2^i}\qquad (i\ge 0).
    \]
    We argue by induction on \(i\). The case \(i=0\) is clear. If
    \[
        L^{(i)}\subseteq L^{2^i},
    \]
    then
    \[
        L^{(i+1)}
        =
        [L^{(i)},L^{(i)}]
        \subseteq
        [L^{2^i},L^{2^i}]
        \subseteq
        L^{2^{i+1}}.
    \]
    Hence the claim follows for all \(i\). If \(L^n=0\) for some \(n\), then for \(i\) sufficiently large one has \(2^i\ge n\), so
    \[
        L^{(i)}\subseteq L^{2^i}=0.
    \]
    Therefore \(L\) is solvable.
\end{proof}

\begin{definition}
    A \emph{solvable ideal} of \(L\) is an ideal which is solvable as a Lie algebra in its own right. Since the sum of two solvable ideals is again solvable, there exists a unique maximal solvable ideal of \(L\), called the \emph{radical} of \(L\), denoted by
    \[
        \operatorname{Rad}(L).
    \]
    The Lie algebra \(L\) is called \emph{semisimple} if
    \[
        \operatorname{Rad}(L)=0.
    \]
\end{definition}

\begin{remark}
    Thus the three notions fit into the chain
    \[
        \text{nilpotent} \;\Longrightarrow\; \text{solvable},
    \]
    while semisimplicity means that \(L\) has no nonzero solvable ideals. In particular, a semisimple Lie algebra cannot contain any nonzero abelian ideal.
\end{remark}

\begin{example}
    Any abelian Lie algebra is nilpotent, hence solvable. On the other hand, the Lie algebra of upper triangular matrices is solvable in general, but need not be nilpotent.
\end{example}

\begin{example}
    Let
    \[
        L=
        \left\{
        \begin{pmatrix}
            a & b \\
            0 & 0
        \end{pmatrix}
        \,\middle|\,
        a,b\in F
        \right\}
        \subseteq \mathfrak{gl}_2(F),
    \]
    Then \(L\) is solvable but not nilpotent.

    Indeed, if
    \[
        x=
        \begin{pmatrix}
            a & b \\
            0 & 0
        \end{pmatrix},
        \qquad
        y=
        \begin{pmatrix}
            c & d \\
            0 & 0
        \end{pmatrix},
    \]
    then
    \[
        [x,y]
        =
        \begin{pmatrix}
            0 & ad-bc \\
            0 & 0
        \end{pmatrix}.
    \]
    Hence
    \[
        [L,L]
        =
        \left\{
        \begin{pmatrix}
            0 & t \\
            0 & 0
        \end{pmatrix}
        \,\middle|\,
        t\in F
        \right\},
        \qquad
        [[L,L],[L,L]]=0.
    \]
    Thus the derived series terminates:
    \[
        L^{(1)}=[L,L],\qquad L^{(2)}=[L^{(1)},L^{(1)}]=0,
    \]
    so \(L\) is solvable.

    On the other hand, \(L\) is not nilpotent, since its descending central series does not reach \(0\). In fact,
    \[
        L^2=[L,L]
        =
        \left\{
        \begin{pmatrix}
            0 & t \\
            0 & 0
        \end{pmatrix}
        \,\middle|\,
        t\in F
        \right\},
    \]
    and for any
    \[
        x=
        \begin{pmatrix}
            a & b \\
            0 & 0
        \end{pmatrix}\in L,\qquad
        y=
        \begin{pmatrix}
            0 & t \\
            0 & 0
        \end{pmatrix}\in L^2,
    \]
    we have
    \[
        [x,y]
        =
        \begin{pmatrix}
            0 & at \\
            0 & 0
        \end{pmatrix}\in L^2.
    \]
    Therefore
    \[
        L^3=[L,L^2]=L^2,
    \]
    and inductively
    \[
        L^n=L^2\neq 0\qquad (n\ge 2).
    \]
    So \(L\) is not nilpotent.
\end{example}

\chapter{Lecture 2 Engel Theorem}

\section{Engel Theorem the first type}

\begin{definition}
    Let \(L\) be a Lie algebra over a field \(F\). An \emph{enveloping algebra} of \(L\) is an associative \(F\)-algebra \(A\) together with a Lie algebra homomorphism
    \[
        \iota:L\to A^{-},
    \]
    where \(A^{-}\) denotes the Lie algebra obtained from \(A\) by equipping the underlying vector space of \(A\) with the bracket
    \[
        [x,y]=xy-yx,
    \]
    such that the associative algebra \(A\) is generated by \(\iota(L)\) (or just writing $A=<L>_{assoc}$ or simply $A=<L>$).

    If, in addition, the pair \((A,\iota)\) satisfies the following universal property: for every associative \(F\)-algebra \(B\) and every Lie algebra homomorphism
    \[
        \varphi:L\to B^{-},
    \]
    there exists a unique associative algebra homomorphism
    \[
        \widetilde{\varphi}:A\to B
    \]
    such that
    \[
        \varphi=\widetilde{\varphi}\circ \iota,
    \]
    then \(A\) is called the \emph{universal enveloping algebra} of \(L\).
\end{definition}

\begin{theorem}
    Let $ A $ be a finite dimensional enveloping algebra of $L$, i.e. $ L\subseteq A^{(-)} $, $ A = <L> $.If $ \exists n \ge 1 $ such that $ \forall a \in L $, $ a^n=0 $. Then $ A $ is nilpotent.
\end{theorem}

\section{Engel Theorem the second type}

In fact we can prove a stronger version.

\begin{definition}
    A subset $S\subseteq L$ is called a \emph{Lie subset} if
    \[
        [a,b]\in S
        \qquad
        \text{for all } a,b\in S.
    \]
    Note that a Lie subset need not be a subspace, and hence need not be a Lie subalgebra.
\end{definition}

\begin{theorem}
    Let $A$ be a finite-dimensional associative algebra over $F$, and let
    \[
        L\subseteq A^{(-)}
    \]
    be a Lie subalgebra such that $A$ is generated by $L$ as an associative algebra. Suppose that
    \[
        L=\mathrm{span}(S)
    \]
    for some Lie subset $S\subseteq L$, and that every element of $S$ is nilpotent in $A$. Then $A$ is a nilpotent associative algebra.
\end{theorem}

\begin{remark}
    In the finite-dimensional situation, the existence of a uniform exponent $n$ is automatic once every element of $S$ is nilpotent in $A$. Indeed, let
    \[
        m=\dim_F A,
    \]
    and let $a\in A$ be nilpotent. Consider the left multiplication operator
    \[
        L_a:A\to A,\qquad x\mapsto ax.
    \]
    Since $a$ is nilpotent, $L_a$ is a nilpotent linear transformation on the $m$-dimensional vector space $A$. Hence
    \[
        L_a^m=0.
    \]
    But
    \[
        L_a^m=L_{a^m},
    \]
    so $a^m=0$. Therefore every nilpotent element of $A$ has nilpotency index at most $\dim_F A$. In particular, if every element of $S$ is nilpotent in $A$, then one may take
    \[
        n=\dim_F A
    \]
    and obtain
    \[
        a^n=0
        \qquad
        \text{for all } a\in S.
    \]
    Thus, in the theorem above, the hypothesis that every element of $S$ is nilpotent in $A$ already implies the existence of a common exponent $n$.
\end{remark}

\begin{theorem}[Zorn's Lemma]
    Let $(P,\leq)$ be a partially ordered set. Assume that every chain in $P$ has an upper bound in $P$. Then $P$ contains a maximal element.
\end{theorem}

\begin{proof}
    We prove the result assuming the Well-Ordering Theorem.

    Suppose, for contradiction, that $P$ has no maximal element. Then for every $x\in P$ there exists some $y\in P$ such that
    \[
        x<y.
    \]
    By the Well-Ordering Theorem, we may choose a well-ordering $\preceq$ on $P$. For each $x\in P$, let $f(x)$ be the $\preceq$-least element of the set
    \[
        \{y\in P: x<y\}.
    \]
    Then
    \[
        x<f(x)\qquad \text{for all }x\in P.
    \]

    We now define, by transfinite recursion, a family $(x_\alpha)$ indexed by ordinals. Suppose that for all $\beta<\alpha$ the elements $x_\beta$ have been defined and form a chain. By hypothesis, the chain
    \[
        C_\alpha:=\{x_\beta:\beta<\alpha\}
    \]
    has an upper bound $u_\alpha\in P$. Define
    \[
        x_\alpha:=f(u_\alpha).
    \]
    Since $u_\alpha$ is an upper bound of $C_\alpha$, for every $\beta<\alpha$ we have
    \[
        x_\beta\leq u_\alpha < x_\alpha.
    \]
    Hence $(x_\alpha)$ is a strictly increasing family:
    \[
        \beta<\alpha \implies x_\beta<x_\alpha.
    \]

    Therefore the map
    \[
        \alpha\longmapsto x_\alpha
    \]
    is injective from the class of all ordinals into the set $P$. This is impossible, since there is no injection from the class of all ordinals into a set.

    This contradiction shows that $P$ must contain a maximal element.
\end{proof}

\begin{lemma}
    Let
    \[
        \mathcal P
        =
        \left\{
        T\subseteq S \;\middle|\;
        T \text{ is a Lie subset and } a^n=0 \text{ for all } a\in T
        \right\},
    \]
    partially ordered by inclusion. Then $\mathcal P$ has a maximal element.
\end{lemma}

\begin{proof}
    We apply Zorn's lemma. Let $\{T_i\}_{i\in I}$ be a chain in $\mathcal P$, and set
    \[
        T=\bigcup_{i\in I}T_i.
    \]
    We claim that $T\in\mathcal P$.

    First, $T$ is a Lie subset. Indeed, if $a,b\in T$, then $a\in T_i$ and $b\in T_j$ for some $i,j\in I$. Since $\{T_i\}_{i\in I}$ is a chain, one of $T_i,T_j$ is contained in the other; say $T_i\subseteq T_j$. Then $a,b\in T_j$, and since $T_j$ is a Lie subset,
    \[
        [a,b]\in T_j\subseteq T.
    \]

    Second, if $a\in T$, then $a\in T_i$ for some $i\in I$, so $a^n=0$ because $T_i\in\mathcal P$.

    Thus $T\in\mathcal P$, and clearly $T$ is an upper bound of the chain. Therefore every chain in $\mathcal P$ has an upper bound. By Zorn's lemma, $\mathcal P$ has a maximal element.
\end{proof}

\begin{proof}[Proof of the theorem]
    Consider the collection
    \[
        \mathcal P=
        \left\{
        T\subseteq S
        \;\middle|\;
        T \text{ is a Lie subset of }L
        \text{ and }
        \langle T\rangle \text{ is a nilpotent associative subalgebra of }A
        \right\},
    \]
    partially ordered by inclusion. Here $\langle T\rangle$ denotes the associative subalgebra of $A$
    generated by $T$.

    We first show that every chain in $\mathcal P$ has an upper bound. Let $\{T_i\}_{i\in I}$ be a chain
    in $\mathcal P$, and set
    \[
        T=\bigcup_{i\in I}T_i.
    \]
    Since the $T_i$ form a chain, $T$ is again a Lie subset of $L$. Let $N=\dim_F A$. Since each
    $\langle T_i\rangle$ is a nilpotent subalgebra of the finite-dimensional algebra $A$, we have
    \[
        \langle T_i\rangle^N=0
        \qquad
        \text{for all }i\in I.
    \]
    Now take any $x_1,\dots,x_N\in \langle T\rangle$. Each $x_j$ belongs to $\langle T_{i_j}\rangle$
    for some $i_j\in I$. Since the family $\{T_i\}$ is totally ordered, there exists $k\in I$ such that
    \[
        T_{i_1}\cup\cdots\cup T_{i_N}\subseteq T_k.
    \]
    Hence
    \[
        x_1,\dots,x_N\in \langle T_k\rangle,
    \]
    and therefore
    \[
        x_1\cdots x_N=0.
    \]
    Thus $\langle T\rangle^N=0$, so $T\in\mathcal P$. By Zorn's lemma, $\mathcal P$ has a maximal
    element, say $S_1$.

    Set
    \[
        B=\langle S_1\rangle.
    \]
    Then $B$ is nilpotent, so there exists $t\ge 1$ such that
    \[
        B^t=0.
    \]

    We claim that if $S_1\ne S$, then there exists $s\in S\setminus S_1$ such that
    \[
        [S_1,s]\subseteq S_1.
    \]
    Assume the contrary. Choose $s_1\in S\setminus S_1$. By assumption, there exists $x_1\in S_1$ such that
    \[
        s_2=[x_1,s_1]\notin S_1.
    \]
    Since $S$ is a Lie subset, $s_2\in S$. Repeating this argument, we construct inductively elements
    \[
        x_1,x_2,\dots \in S_1,
        \qquad
        s_1,s_2,\dots \in S\setminus S_1
    \]
    such that
    \[
        s_{r+1}=[x_r,s_r]
        \qquad
        (r\ge 1).
    \]
    Thus
    \[
        s_{r+1}=\operatorname{ad}(x_r)\cdots \operatorname{ad}(x_1)(s_1).
    \]

    We now show that this process cannot continue for $2t$ steps. Indeed, for any $y_1,\dots,y_r\in B$,
    the iterated commutator
    \[
        \operatorname{ad}(y_r)\cdots \operatorname{ad}(y_1)(s_1)
    \]
    is a sum of terms of the form
    \[
        \pm u\, s_1\, v,
    \]
    where $u,v\in B$ and the total number of factors from $B$ occurring in $u$ and $v$ is $r$.
    Therefore, if $r\ge 2t$, then in each such term either $u\in B^t$ or $v\in B^t$, so the term is $0$.
    Hence
    \[
        \operatorname{ad}(y_r)\cdots \operatorname{ad}(y_1)(s_1)=0
        \qquad
        \text{for all } r\ge 2t.
    \]
    Applying this with $y_i=x_i\in S_1\subseteq B$, we obtain
    \[
        s_{2t+1}=0,
    \]
    which is impossible since $0\in S_1$. This contradiction proves the claim.

    Choose $s\in S\setminus S_1$ such that
    \[
        [S_1,s]\subseteq S_1.
    \]
    Let
    \[
        S_2=S_1\cup\{s\}.
    \]
    Then $S_2$ is a Lie subset of $L$: indeed, brackets of elements of $S_1$ remain in $S_1$,
    \([S_1,s]\subseteq S_1\), \([s,S_1]\subseteq S_1\) by antisymmetry, and \([s,s]=0\in S_1\).

    We next show that $\langle S_2\rangle$ is nilpotent. Since $[S_1,s]\subseteq S_1$, we have
    \[
        [B,s]\subseteq B.
    \]
    Hence for every $b\in B$,
    \[
        bs=sb+[b,s],
        \qquad
        [b,s]\in B.
    \]
    By repeated use of this identity, every product in the algebra generated by $B$ and $s$ can be written
    as a linear combination of terms of the form
    \[
        s^r b,
        \qquad
        b\in B,
    \]
    with $r\ge 0$. Since $s^m=0$, we may assume $0\le r\le m-1$.

    Now consider a product of length $mt$ in $\langle S_2\rangle$. After rewriting it as above, each term
    has the form $s^r b$ with $0\le r\le m-1$. Since the original length is $mt$, such a term must contain
    at least
    \[
        mt-r\ge mt-(m-1)\ge t
    \]
    factors coming from $B$. Therefore $b\in B^t=0$, and every such term is zero. It follows that
    \[
        \langle S_2\rangle^{\,mt}=0.
    \]
    Thus $\langle S_2\rangle$ is nilpotent.

    But $S_1\subsetneq S_2$ and $S_2\in\mathcal P$, contradicting the maximality of $S_1$. Hence
    $S_1=S$. Therefore
    \[
        A=\langle L\rangle=\langle S\rangle=\langle S_1\rangle
    \]
    is nilpotent.
\end{proof}

\section{Engel Theorem the third type}

\begin{definition}
    Let $L$ be a Lie algebra over a field $F$. An $L$-module is a vector space $V$ over $F$ together with a bilinear map
    \[
        L\times V \to V,\qquad (x,v)\mapsto x\cdot v,
    \]
    such that
    \[
        [x,y]\cdot v = x\cdot (y\cdot v)-y\cdot (x\cdot v)
    \]
    for all $x,y\in L$ and $v\in V$. In other words, $V$ is an $L$-module if and only if there exists a Lie algebra homomorphism
    \[
        \phi:L\to \mathfrak{gl}(V).
    \]
    such that
    \[
        \phi(x)(v) = x\cdot v
    \]
    for all $x\in L$ and $v\in V$.
\end{definition}

\begin{proposition}
    Let $L$ be a Lie algebra over a field $F$, and let $M$ be a vector space over $F$. Suppose a bilinear operation
    \[
        L\times M\to M,\qquad (a,m)\mapsto a\cdot m
    \]
    is given. Then the following conditions are equivalent:

    \begin{enumerate}
        \item The vector space $L\oplus M$ becomes a Lie algebra under the bracket defined by
              \[
                  [a+m,\ b+n]=[a,b]+a\cdot n-b\cdot m
                  \qquad (a,b\in L,\ m,n\in M),
              \]
              equivalently,
              \[
                  [a,m]=a\cdot m,\qquad [m,a]=-\,a\cdot m
              \]
              for all $a\in L$, $m\in M$, and
              \[
                  [m,n]=0
              \]
              for all $m,n\in M$.

        \item For all $a,b\in L$ and $m\in M$,
              \[
                  [a,b]\cdot m=a\cdot (b\cdot m)-b\cdot (a\cdot m).
              \]

        \item The map
              \[
                  \rho:L\to \mathfrak{gl}_F(M),\qquad \rho(a)(m)=a\cdot m,
              \]
              is a Lie algebra homomorphism.
    \end{enumerate}

    When these equivalent conditions hold, we say that $M$ is an $L$-module.
\end{proposition}

\begin{proof}
    We prove $(1)\Rightarrow(2)\Rightarrow(3)\Rightarrow(1)$.

    \medskip

    $(1)\Rightarrow(2)$. Assume that $L\oplus M$ is a Lie algebra with the stated bracket. Since every Lie algebra satisfies the Jacobi identity, for any $a,b\in L$ and $m\in M$ we have
    \[
        [a,[b,m]]+[b,[m,a]]+[m,[a,b]]=0.
    \]
    Now
    \[
        [b,m]=b\cdot m,\qquad [m,a]=-\,a\cdot m,\qquad [a,b]\in L.
    \]
    Because $M$ is abelian, this gives
    \[
        [a,b\cdot m]+[b,-\,a\cdot m]+[m,[a,b]]=0.
    \]
    Using again the definition of the bracket,
    \[
        a\cdot (b\cdot m)-b\cdot (a\cdot m)-[a,b]\cdot m=0.
    \]
    Hence
    \[
        [a,b]\cdot m=a\cdot (b\cdot m)-b\cdot (a\cdot m).
    \]

    \medskip

    $(2)\Rightarrow(3)$. Define
    \[
        \rho:L\to \mathfrak{gl}_F(M),\qquad \rho(a)(m)=a\cdot m.
    \]
    Since the action is bilinear, $\rho$ is linear. For any $a,b\in L$ and $m\in M$, condition (2) gives
    \[
        \rho([a,b])(m)=[a,b]\cdot m
        =a\cdot (b\cdot m)-b\cdot (a\cdot m)
        =\rho(a)\rho(b)(m)-\rho(b)\rho(a)(m).
    \]
    Therefore
    \[
        \rho([a,b])=[\rho(a),\rho(b)],
    \]
    so $\rho$ is a Lie algebra homomorphism.

    \medskip

    $(3)\Rightarrow(1)$. Assume that $\rho:L\to \mathfrak{gl}_F(M)$ is a Lie algebra homomorphism, and define a bracket on $L\oplus M$ by
    \[
        [a+m,\ b+n]=[a,b]+\rho(a)(n)-\rho(b)(m).
    \]
    This bracket is bilinear and skew-symmetric. It remains to verify the Jacobi identity.

    Take $a,b,c\in L$ and $m,n,p\in M$. Write
    \[
        x=a+m,\qquad y=b+n,\qquad z=c+p.
    \]
    A direct calculation shows that the $L$-component of
    \[
        [x,[y,z]]+[y,[z,x]]+[z,[x,y]]
    \]
    is
    \[
        [a,[b,c]]+[b,[c,a]]+[c,[a,b]],
    \]
    which is zero because $L$ is a Lie algebra.

    For the $M$-component, expanding the definition of the bracket yields
    \[
        \rho(a)\rho(b)(p)-\rho(a)\rho(c)(n)-\rho([b,c])(m)
        +\rho(b)\rho(c)(m)-\rho(b)\rho(a)(p)-\rho([c,a])(n)
    \]
    \[
        +\rho(c)\rho(a)(n)-\rho(c)\rho(b)(m)-\rho([a,b])(p).
    \]
    Grouping the terms by $m,n,p$, we obtain
    \[
        \bigl(\rho(b)\rho(c)-\rho(c)\rho(b)-\rho([b,c])\bigr)(m)
    \]
    \[
        +\bigl(\rho(c)\rho(a)-\rho(a)\rho(c)-\rho([c,a])\bigr)(n)
    \]
    \[
        +\bigl(\rho(a)\rho(b)-\rho(b)\rho(a)-\rho([a,b])\bigr)(p).
    \]
    Since $\rho$ is a Lie algebra homomorphism, each bracketed expression is zero. Hence the $M$-component also vanishes, so the Jacobi identity holds.

    Therefore $L\oplus M$ is a Lie algebra with the stated bracket.

    This proves that the three conditions are equivalent.
\end{proof}

\begin{theorem}
    Let $L$ be a finite dimensional Lie algebra over a field $F$. If $\forall a \in L$, $\operatorname{ad}(a) $ is nilpotent, then $L$ is nilpotent.
\end{theorem}

\begin{proof}
    Note that since $L$ is finite-dimensional, say $\dim L=n<\infty$, and for every $a\in L$ we have
    \[
        (\operatorname{ad} a)^n=0.
    \]
    Hence we may apply the previous theorem: if for all $a\in L$,
    \[
        (\operatorname{ad} a)^n=0,
    \]
    then
    \[
        (\operatorname{ad} L)^{\,n}=0,
    \]
    and therefore
    \[
        L^{n+1}=0.
    \]
    So $L$ is nilpotent.
\end{proof}

\begin{remark}
    If $L$ is infinite-dimensional, one should distinguish two cases:
    \[
        \left\{
        \begin{array}{l}
            \operatorname{ad} a \text{ nilpotent with a uniform bound } n
            \ \Longrightarrow\
            L \text{ nilpotent}, \\[4pt]
            \text{no uniform bound on the nilpotency indices}
            \ \nRightarrow\
            L \text{ nilpotent}.
        \end{array}
        \right.
    \]
\end{remark}

\begin{remark}
    There are three related versions, with
    \[
        \text{type }(2)\Longrightarrow \text{type }(1)\Longrightarrow \text{type }(3).
    \]

    Type (2): Let $S\subseteq L$ be a Lie subset. And $S$ generates $L$ as a Lie algebra. Assume there exists $n$ such that for every $a\in S$
    \[
        a^n =0.
    \]
    Then $A=<L>$ is a nilpotent algebra. (When $A$ is a finite-dimensional associative algebra).

    In particular, if we choose $S$ to be $L$, then we get type (1).

    Type (1): Let $ A $ be a finite dimensional enveloping algebra of $L$, i.e. $ L\subseteq A^{(-)} $, $ A = <L> $.If $ \exists n \ge 1 $ such that $ \forall a \in L $, $ a^n=0 $. Then $ A $ is nilpotent.

    And consider $A$ as $\operatorname{ad}(L)$, we get type (3).

    Type (3):    Let $L$ be a finite dimensional Lie algebra over a field $F$. If $\forall a \in L$, $\operatorname{ad}(a) $ is nilpotent, then $L$ is nilpotent.
\end{remark}

We also have the following truth:

\begin{corollary}
    \ Let $M$ be an $L$-module with $\dim M<\infty$, that is, let
    \[
        \varphi:L\to \mathfrak{gl}_F(M)
    \]
    be a representation. Assume that for every $a\in L$, the operator $\varphi(a)$ is nilpotent. Equivalently, for every $a\in L$ there exists $n$ such that
    \[
        \underbrace{a\bigl(a(\cdots(a}_{n\text{ times}}\,M)\cdots)\bigr)=0.
    \]
    Then there exists $N$ such that
    \[
        \underbrace{L\cdots L}_{N\text{ times}}M=0.
    \]
    In other words, if every $\varphi(a)$ is nilpotent, then the Lie algebra $\varphi(L)$ is nilpotent. Hence there exists a basis of $M$ in which every element of $\varphi(L)$ is represented by a strictly upper triangular matrix:
    \[
        \varphi(L)\subseteq
        \left\{
        \begin{pmatrix}
            0      & *      & \cdots & *      \\
            0      & 0      & \cdots & *      \\
            \vdots & \vdots & \ddots & \vdots \\
            0      & 0      & \cdots & 0
        \end{pmatrix}
        \right\}.
    \]
\end{corollary}

\begin{remark}
    This corollary is the module-theoretic form of Engel's theorem. Indeed, the representation
    \[
        \varphi:L\to \mathfrak{gl}_F(M)
    \]
    has image $\varphi(L)$, which is a Lie subalgebra of $\mathfrak{gl}_F(M)$. By assumption, every element of $\varphi(L)$ is a nilpotent endomorphism of the finite-dimensional vector space $M$. Hence, by Engel's theorem, the Lie algebra $\varphi(L)$ is nilpotent.

    The conclusion
    \[
        \underbrace{L\cdots L}_{N\text{ times}}M=0
    \]
    means that there exists a uniform integer $N$ such that for all $a_1,\dots,a_N\in L$ and all $m\in M$,
    \[
        a_1\bigl(a_2(\cdots(a_Nm)\cdots)\bigr)=0.
    \]
    Thus the action is nilpotent not only for each fixed $a\in L$, but uniformly for all sufficiently long iterated actions.

    A further consequence of Engel's theorem is that there exists a basis of $M$ with respect to which every operator $\varphi(a)$ is represented by a strictly upper triangular matrix. Equivalently, $M$ admits a chain of subspaces
    \[
        0=M_0\subset M_1\subset \cdots \subset M_r=M,
        \qquad \dim M_i=i,
    \]
    such that
    \[
        L\cdot M_i\subseteq M_{i-1}
        \qquad (1\le i\le r).
    \]

    Explanation: By Engel's theorem, $\varphi(L)$ is nilpotent. Hence
    \[
        \underbrace{L\cdots L}_{N\text{ times}}M=0
    \]
    Take the minimal $N$ such that $L^N M=0$ and $L^{N-1}M \neq 0$. Then we have a non-zero vector $v_1$ in $L^{N-1}M$ which is killed by $\varphi(L)$. Denote $M_1 = \operatorname{span}(v_1)$. Then $L\cdot M_1 \subseteq M_0$. Now consider $M/M_1$. By Engel's theorem and by induction hypothesis, $\varphi(L)$ is nilpotent on $M/M_1$. Hence there exists a basis of $M/M_1$ with respect to which every element of $\varphi(L)$ is represented by a strictly upper triangular matrix. Equivalently, $M/M_1$ admits a chain of subspaces
    \[
        0=M_1/M_1\subset M_2/M_1\subset \cdots \subset M_r/M_1=M/M_1,
        \qquad \dim M_i/M_1=i,
    \]
    such that
    \[
        L\cdot (M_i/M_1)\subseteq M_{i-1}/M_1
        \qquad (1\le i\le r).
    \]
    Now we can take a basis of $M_1$ and $M/M_1$ to get a basis of $M$ with respect to which every element of $\varphi(L)$ is represented by a strictly upper triangular matrix.

    In particular, this corollary reduces to the usual form of Engel's theorem when one takes $M=L$ and $\varphi=\operatorname{ad}$.
\end{remark}

\begin{remark}
    According to this corollary and type (2), we can rewrite it as:

    Let $M$ be an $L$-module with $\dim M<\infty$, that is, let
    \[
        \varphi:L\to \mathfrak{gl}_F(M)
    \]
    be a representation. Assume that for every $a\in S$, where $S$ is a Lie subset of $L$ and $<S>=L$, the operator $\varphi(a)$ is nilpotent. Equivalently, for every $a\in S$ there exists $n$ such that
    \[
        \underbrace{a\bigl(a(\cdots(a}_{n\text{ times}}\,M)\cdots)\bigr)=0.
    \]
    Then there exists $N$ such that
    \[
        \underbrace{L\cdots L}_{N\text{ times}}M=0.
    \]
    For this result, we do not need $L$ or the enveloping algebra $A=<L>$ to be finite-dimensional now, we only need such $n$ is a universal bound.
\end{remark}

\chapter{Lecture 3: Lie's theorem and Jordan-Chevalley decomposition}

\section{Lie's theorem}

\begin{lemma}
    Let $F$ be a field with $\operatorname{char} F=0$ or $\operatorname{char} F>n$.
    Let $a,b\in M_n(F)$ satisfy
    \[
        [a,[a,b]]=0.
    \]
    Then $[a,b]$ is nilpotent. In fact,
    \[
        [a,b]^n=0.
    \]
\end{lemma}

\begin{proof}
    Set
    \[
        \delta:=\operatorname{ad}(a),\qquad \delta(x)=[a,x].
    \]
    Since $\delta$ is an inner derivation of the associative algebra $M_n(F)$, we have
    \[
        \delta(xy)=\delta(x)y+x\delta(y)
        \qquad\text{for all }x,y\in M_n(F).
    \]
    Our assumption is precisely
    \[
        \delta^2(b)=0.
    \]

    Now let us prove the following two facts:

    \medskip
    \noindent
    \textit{Claim 1.} For every $m\ge 1$ and every $r>m$,
    \[\operatorname{End}
        \delta^r(b^m)=0.
    \]

    \noindent
    Indeed, each application of $\delta$ to the product $b^m$ differentiates one of the $m$
    factors. Since $\delta^2(b)=0$, no factor can be differentiated twice. Hence after more
    than $m$ applications, every term vanishes.

    \medskip
    \noindent
    \textit{Claim 2.} For every $m\ge 1$,
    \[
        \delta^m(b^m)=m!\,\delta(b)^m=m!\,[a,b]^m.
    \]

    \noindent
    We prove this by induction on $m$. The case $m=1$ is clear. Assume the formula holds
    for $m-1$. Using the derivation rule,
    \[
        \delta(b^m)=\delta(b)b^{m-1}+b\,\delta(b^{m-1}).
    \]
    Applying $\delta^{m-1}$ and using Leibniz repeatedly, together with $\delta^2(b)=0$,
    one finds that the only surviving terms are those in which each of the $m$ copies of $b$
    is hit exactly once. There are exactly $m!$ such terms, and each contributes
    \[
        \delta(b)^m.
    \]
    Hence
    \[
        \delta^m(b^m)=m!\,\delta(b)^m.
    \]

    \medskip
    Now apply Cayley--Hamilton to $b$:
    \[
        b^n+\beta_{n-1}b^{n-1}+\cdots+\beta_1 b+\beta_0 I=0.
    \]
    Apply $\delta^n$ to both sides. By Claim 1, for every $i<n$ we have
    \[
        \delta^n(b^i)=0,
    \]
    and also $\delta^n(I)=0$. Therefore
    \[
        \delta^n(b^n)=0.
    \]
    By Claim 2 with $m=n$, this becomes
    \[
        n!\,[a,b]^n=0.
    \]
    Since $\operatorname{char} F=0$ or $\operatorname{char} F>n$, we have $n!\neq 0$ in $F$.
    Hence
    \[
        [a,b]^n=0.
    \]
    Therefore $[a,b]$ is nilpotent.
\end{proof}

\begin{theorem}[Lie's theorem]
    Assume that $\operatorname{char} F = 0$ or $\operatorname{char} F > \dim_F M$. $F$ is algebraically closed.
    Let $L$ be a solvable Lie algebra and let $M$ be an $L$-module with
    \[
        \dim_F M = n < \infty.
    \]
    Then there exists a basis of $M$ such that, relative to this basis, the image of $L$ (from here when talk about $L$ acting on its module $M$, we simply use $L$ instead of $\varphi(L)$)
    consists of upper triangular matrices; that is,
    \[
        L \subseteq
        \left\{
        \begin{pmatrix}
            *      & *      & \cdots & *      \\
            0      & *      & \cdots & *      \\
            \vdots & \ddots & \ddots & \vdots \\
            0      & \cdots & 0      & *
        \end{pmatrix}
        \right\}.
    \]
\end{theorem}

\begin{proof}
    We argue first that $M$ has a common eigenvector for the action of $L$.

    Assume first that $M$ is irreducible as an $L$-module. Since we are mainly concerned
    with the image of
    \[
        L \to \mathfrak{gl}_F(M),
    \]
    we may replace $L$ by its image and assume that $L \subseteq \mathfrak{gl}_F(M)$.

    Because $L$ is solvable, its derived series
    \[
        L \supseteq L^{[1]} \supseteq L^{[2]} \supseteq \cdots
    \]
    terminates at $0$. Hence there exists an integer $s \ge 0$ such that
    \[
        L^{[s]} \neq 0,
        \qquad
        L^{[s+1]} = [L^{[s]},L^{[s]}]=0.
    \]
    If $s=0$, then $L$ is abelian, and since $F$ is algebraically closed, there exists
    an eigenvector $v \in M$ for some $x \in L$. As $L$ is abelian, every element of $L$
    preserves the eigenspace of $x$, and by irreducibility this gives a common eigenvector
    for all $x \in L$.

    Now assume $s \ge 1$. Choose $a \in L^{[s]}$ and $b \in L$. Then
    \[
        [a,[a,b]] \in [L^{[s]},L^{[s]}] = L^{[s+1]} = 0,
    \]
    so
    \[
        (\operatorname{ad} a)^2(b)=0.
    \]
    Consider the set
    \[
        S=\{[a,b]\mid a\in L^{[s]},\ b\in L\}.
    \]
    This is a Lie subset of $L$, and every element of $S$ is nilpotent in $\mathfrak{gl}_F(M)$.
    By Engel's theorem,
    \[
        \operatorname{span}_F(S)=[L^{[s]},L]
    \]
    is a nilpotent Lie algebra of linear transformations on $M$.

    Now let $I$ be any ideal of $L$. Then $IM$ is an $L$-submodule of $M$, because for
    $x\in L$, $y\in I$, and $m\in M$,
    \[
        x\cdot (y\cdot m)=y\cdot (x\cdot m)+[x,y]\cdot m \in IM,
    \]
    since $[x,y]\in I$. As $M$ is irreducible, it follows that
    \[
        IM=0 \quad \text{or} \quad IM=M.
    \]
    Applying this to the ideal $I=L^{[s]}$, we cannot have $L^{[s]}M=M$, because
    $L^{[s]}$ acts nilpotently on $M$. Hence
    \[
        L^{[s]}M=0,
    \]
    that is, $L^{[s]}$ acts trivially on $M$.

    Thus every element of $L^{[s]}$ commutes with the action of $L$ on $M$. In particular,
    for any $c\in L^{[s-1]}$ and $d\in L^{[s-1]}$,
    \[
        [c,d]\in L^{[s]}
    \]
    acts as $0$ on $M$. Therefore $L^{[s]}=0$ in the image of $L$ in $\mathfrak{gl}_F(M)$,
    contrary to the choice of $s$. This contradiction shows that $s=0$, so $L$ is abelian
    (on $M$), and hence $M$ contains a common eigenvector for all elements of $L$.

    Now let $0\neq v_1\in M$ be such a common eigenvector. Then $Fv_1$ is a one-dimensional
    $L$-submodule of $M$. Consider the quotient $M/Fv_1$, which is again a finite-dimensional
    $L$-module. By induction on $\dim_F M$, there exists a basis of $M/Fv_1$ with respect to
    which the induced action of $L$ is upper triangular. Lifting this basis to $M$ and adjoining
    $v_1$, we obtain a basis of $M$ relative to which every element of $L$ is represented by an
    upper triangular matrix.

    Therefore there exists a basis of $M$ such that the image of $L$ consists of upper triangular
    matrices.
\end{proof}

\begin{proposition}
    Let $L$ be a solvable Lie algebra over $F$ with $\operatorname{char} F = 0$ or $\operatorname{char} F > \dim_F M$. $F$ is algebraically closed. Then the followings are equivalent:
    \begin{enumerate}
        \item Any irreducible finite dimensional $L$-module is 1-dimensional.
        \item If $M$ is a finite dimensional $L$-module, then there exist a basis of $M$ such that the image of $L$ consists of upper triangular matrices.
        \item If $M$ is any finite dimensional $L$-module. Then $[L,L]$ acts nilpotently on $M$.
    \end{enumerate}
\end{proposition}

\begin{proof}
    We prove
    \[
        (2)\Rightarrow(1),\qquad (1)\Rightarrow(2),\qquad (2)\Rightarrow(3),\qquad (3)\Rightarrow(1).
    \]

    $(2)\Rightarrow(1)$.
    Let $M$ be an irreducible finite-dimensional $L$-module. By (2), there exists a basis of $M$ relative to which every operator in the image of $L$ is upper triangular. Then the subspace spanned by the first basis vector is stable under $L$, hence is a nonzero submodule of $M$. Since $M$ is irreducible, this submodule must be all of $M$. Therefore $\dim_F M=1$.

    $(1)\Rightarrow(2)$.
    We argue by induction on $\dim_F M$.

    If $\dim_F M=1$, there is nothing to prove. Assume $\dim_F M>1$, and suppose the statement is known for all modules of smaller dimension. Since $M$ is finite-dimensional, it contains an irreducible submodule $N$. By (1), we have $\dim_F N=1$.

    Choose $0\neq v_1\in N$, so that $N=Fv_1$. Then $N$ is an $L$-submodule, and hence $M/N$ is again a finite-dimensional $L$-module. By the induction hypothesis, there exists a basis
    \[
        \overline{v_2},\dots,\overline{v_n}
    \]
    of $M/N$ such that the induced action of $L$ on $M/N$ is given by upper triangular matrices. Choose representatives $v_2,\dots,v_n\in M$ of these cosets. Then
    \[
        v_1,v_2,\dots,v_n
    \]
    is a basis of $M$.

    Since $Fv_1$ is $L$-stable, every element of $L$ sends $v_1$ into $Fv_1$. Since the induced action on $M/Fv_1$ is upper triangular with respect to
    \[
        \overline{v_2},\dots,\overline{v_n},
    \]
    it follows that the action of $L$ on $M$ is upper triangular with respect to
    \[
        v_1,v_2,\dots,v_n.
    \]
    Thus (2) holds.

    $(2)\Rightarrow(3)$.
    Let $M$ be a finite-dimensional $L$-module, and choose a basis of $M$ such that every element of the image of $L$ is represented by an upper triangular matrix. Then for any $x,y\in L$,
    \[
        \rho([x,y])=[\rho(x),\rho(y)].
    \]
    Since the commutator of two upper triangular matrices has zero diagonal entries, $\rho([x,y])$ is strictly upper triangular. Hence every element of $\rho([L,L])$ is strictly upper triangular, and therefore nilpotent. Thus $[L,L]$ acts nilpotently on $M$.

    $(3)\Rightarrow(1)$.
    Let $M$ be an irreducible finite-dimensional $L$-module. By (3), every element of $[L,L]$ acts nilpotently on $M$. Hence the Lie algebra $\rho([L,L])\subseteq \mathfrak{gl}(M)$ consists entirely of nilpotent endomorphisms. By Engel's theorem, there exists a nonzero vector $v\in M$ such that
    \[
        \rho([L,L])v=0.
    \]
    Thus $[L,L]$ annihilates $v$.

    Let $N=U(L)v$ be the submodule generated by $v$. Since $[L,L]v=0$, the action of $L$ on $N$ factors through the abelian Lie algebra $L/[L,L]$. Therefore $N$ is a finite-dimensional module for the abelian Lie algebra $L/[L,L]$.

    Because $F$ is algebraically closed, every irreducible finite-dimensional module over an abelian Lie algebra is $1$-dimensional. Hence $N$ contains a $1$-dimensional submodule. Since $M$ is irreducible and $0\neq N\subseteq M$, we have $N=M$. Therefore $M$ itself contains a $1$-dimensional submodule, and irreducibility forces $\dim_F M=1$.

    This proves (1).

    \smallskip

    Therefore
    \[
        (1)\Longleftrightarrow(2)\Longleftrightarrow(3).
    \]
\end{proof}

\begin{remark}
    The assumption that $F$ is algebraically closed is used in the last step of the proof of $(3)\Rightarrow(1)$, namely in the fact that every irreducible finite-dimensional module over the abelian Lie algebra $L/[L,L]$ is $1$-dimensional. Equivalently, one needs the existence of eigenvalues and eigenvectors for commuting linear operators.
\end{remark}

\begin{example}
    Assume that $\operatorname{char} F = 2$. Then $\mathfrak{gl}(2,F)$ is solvable.

    Indeed, every element of $\mathfrak{gl}(2,F)$ has the form
    \[
        \begin{pmatrix}
            \alpha & \gamma \\
            \beta  & \delta
        \end{pmatrix},
        \qquad \alpha,\beta,\gamma,\delta \in F.
    \]
    A direct computation shows that the commutator of two such matrices is always of the form
    \[
        \left[
            \begin{pmatrix}
                \alpha & \gamma \\
                \beta  & \delta
            \end{pmatrix},
            \begin{pmatrix}
                \alpha' & \gamma' \\
                \beta'  & \delta'
            \end{pmatrix}
            \right]
        =
        \begin{pmatrix}
            d & x \\
            y & d
        \end{pmatrix},
    \]
    for some $d,x,y \in F$.
    Hence
    \[
        \mathfrak{gl}(2,F)^{(1)}=[\mathfrak{gl}(2,F),\mathfrak{gl}(2,F)]
        \subseteq
        \left\{
        \begin{pmatrix}
            d & x \\
            y & d
        \end{pmatrix}
        :\ d,x,y\in F
        \right\}.
    \]

    Now if
    \[
        A
        =
        \begin{pmatrix}
            \alpha & \gamma \\
            \beta  & \alpha
        \end{pmatrix}
        \qquad (\text{since } \operatorname{char} F=2),
    \]
    and similarly
    \[
        B=
        \begin{pmatrix}
            \alpha' & \gamma' \\
            \beta'  & \alpha'
        \end{pmatrix},
    \]
    then their commutator is of the form
    \[
        [A,B]=
        \begin{pmatrix}
            d & 0 \\
            0 & d
        \end{pmatrix}
    \]
    for some $d\in F$.
    Therefore
    \[
        \mathfrak{gl}(2,F)^{(2)}
        \subseteq
        \left\{
        \begin{pmatrix}
            d & 0 \\
            0 & d
        \end{pmatrix}
        :\ d\in F
        \right\},
    \]
    and so
    \[
        \mathfrak{gl}(2,F)^{(3)}=0.
    \]
    Thus $\mathfrak{gl}(2,F)$ is solvable.
\end{example}

\section{Generalized eigenvectors and generalized eigenspaces}

\begin{definition}
    Let $V$ be a finite-dimensional vector space over a field $F$ where $F$ is algebraically closed, and let
    \[
        T\in \operatorname{End}_F(V).
    \]
    Suppose that the characteristic polynomial of $T$ splits as
    \[
        f(x)=(x-\lambda_1)^{a_1}\cdots (x-\lambda_k)^{a_k},
    \]
    where $\lambda_1,\dots,\lambda_k$ are distinct. For each $i$, define
    \[
        V_{\lambda_i}:=\ker (T-\lambda_i I)^{a_i}.
    \]
    This is called the generalized eigenspace of $T$ corresponding to $\lambda_i$.
\end{definition}

\begin{theorem}
    With the notation above,
    \[
        V=V_{\lambda_1}\oplus \cdots \oplus V_{\lambda_k}.
    \]
\end{theorem}

\begin{proof}
    For each $i$, let
    \[
        P_i(x)=\frac{f(x)}{(x-\lambda_i)^{a_i}}
        =\prod_{j\ne i}(x-\lambda_j)^{a_j}.
    \]
    Since the polynomials $P_1(x),\dots,P_k(x)$ have greatest common divisor $1$, there exist polynomials
    \[
        Q_1(x),\dots,Q_k(x)\in F[x]
    \]
    such that
    \[
        Q_1(x)P_1(x)+\cdots+Q_k(x)P_k(x)=1.
    \]
    Replacing $x$ by $T$, we obtain
    \[
        Q_1(T)P_1(T)+\cdots+Q_k(T)P_k(T)=I_V.
    \]
    Hence every $v\in V$ can be written as
    \[
        v=\sum_{i=1}^k Q_i(T)P_i(T)v.
    \]
    Now
    \[
        (T-\lambda_i I)^{a_i}Q_i(T)P_i(T)
        =Q_i(T)f(T)=0,
    \]
    because $f(T)=0$ by the Cayley--Hamilton theorem. Therefore
    \[
        Q_i(T)P_i(T)v\in \ker (T-\lambda_i I)^{a_i}=V_{\lambda_i},
    \]
    and so
    \[
        V=V_{\lambda_1}+\cdots+V_{\lambda_k}.
    \]

    To show that the sum is direct, let
    \[
        v_1+\cdots+v_k=0,
        \qquad v_i\in V_{\lambda_i}.
    \]
    Fix $i$. Since $P_i(\lambda_i)\ne 0$, the polynomial $P_i(x)$ is relatively prime to $(x-\lambda_i)^{a_i}$. Hence there exist polynomials $A_i(x),B_i(x)\in F[x]$ such that
    \[
        A_i(x)P_i(x)+B_i(x)(x-\lambda_i)^{a_i}=1.
    \]
    Replacing $x$ by $T$ and applying to $v_i\in V_{\lambda_i}$, we get
    \[
        A_i(T)P_i(T)v_i=v_i,
    \]
    because $(T-\lambda_i I)^{a_i}v_i=0$. Thus $P_i(T)$ is invertible on $V_{\lambda_i}$.

    On the other hand, if $j\ne i$, then
    \[
        P_i(T)v_j=0,
    \]
    because $P_i(x)$ contains the factor $(x-\lambda_j)^{a_j}$ and
    \[
        (T-\lambda_j I)^{a_j}v_j=0.
    \]
    Applying $P_i(T)$ to the relation
    \[
        v_1+\cdots+v_k=0,
    \]
    we obtain
    \[
        P_i(T)v_i=0.
    \]
    Since $P_i(T)$ is invertible on $V_{\lambda_i}$, it follows that $v_i=0$. As $i$ was arbitrary, all $v_i$ are zero. Therefore the sum is direct.
\end{proof}

\chapter{Homework}

\section{Homework week 1}

\begin{exercise}
    Prove the simplicity of Lie algebras $\mathfrak{sl}_n$, $\mathfrak{o}_n$, $\mathfrak{sp}_{2n}$ and $W(n)$.
\end{exercise}

I can only prove when base field is of characteristic 0 and $\neq 2$. (Or $\operatorname{Char}F \nmid n$, which is essentially the same. If $\operatorname{Char}F \mid n$ or $\operatorname{Char}F = 2$, then either $I$ will contained in the center of $L$, or the alternativity become symmetry, things could be more complicated).

\textbf{The simplicity of $\mathfrak{sl}_n$}

\begin{proof}
    Let $I\neq 0$ be an ideal of $\mathfrak{sl}_n(F)$. We prove that $I=\mathfrak{sl}_n(F)$.

    Write
    \[
    \mathfrak{sl}_n(F)=\mathfrak{h}\oplus \bigoplus_{i\neq j}Fe_{ij},
    \]
    where $\mathfrak{h}$ is the subalgebra of diagonal trace-zero matrices. Choose $0\neq x\in I$ and write
    \[
    x=y+\sum_{t=1}^s c_t e_{i_tj_t},
    \qquad y\in\mathfrak{h},
    \]
    where the pairs $(i_t,j_t)$ are distinct and $c_t\neq 0$ for all $t$.

    We first show that $I$ contains some off-diagonal matrix unit $e_{kl}$.

    If $s=0$, then $x=y\in\mathfrak{h}$ is non-zero. Writing
    \[
    y=\operatorname{diag}(d_1,\dots,d_n),
    \]
    we can choose $k\neq l$ such that $d_k\neq d_l$. Then
    \[
    [y,e_{kl}]=(d_k-d_l)e_{kl}\in I,
    \]
    hence $e_{kl}\in I$.

    Assume now that $s\geq 1$. For
    \[
    h=\operatorname{diag}(\lambda_1,\dots,\lambda_n)\in\mathfrak{h},
    \]
    define
    \[
    \alpha_{ij}(h)=\lambda_i-\lambda_j.
    \]
    Since $\operatorname{char}F=0$, the field $F$ is infinite, so we may choose $h\in\mathfrak{h}$ such that
    \[
    \mu_t:=\alpha_{i_tj_t}(h)\neq 0
    \qquad\text{and}\qquad
    \mu_1,\dots,\mu_s
    \text{ are pairwise distinct.}
    \]
    Because $\mathfrak{h}$ is abelian, for every $r\geq 1$,
    \[
    (\operatorname{ad}h)^r(x)=\sum_{t=1}^s c_t\mu_t^r e_{i_tj_t}.
    \]
    Put
    \[
    z_r:=(\operatorname{ad}h)^r(x)\in I,
    \qquad r=1,\dots,s.
    \]
    Then
    \[
    \begin{pmatrix}
    z_1\\
    z_2\\
    \vdots\\
    z_s
    \end{pmatrix}
    =
    \begin{pmatrix}
    \mu_1 & \mu_2 & \cdots & \mu_s\\
    \mu_1^2 & \mu_2^2 & \cdots & \mu_s^2\\
    \vdots & \vdots & \ddots & \vdots\\
    \mu_1^s & \mu_2^s & \cdots & \mu_s^s
    \end{pmatrix}
    \begin{pmatrix}
    c_1e_{i_1j_1}\\
    c_2e_{i_2j_2}\\
    \vdots\\
    c_se_{i_sj_s}
    \end{pmatrix}.
    \]
    The coefficient matrix is invertible: it is the product of $\operatorname{diag}(\mu_1,\dots,\mu_s)$ and a Vandermonde matrix in $\mu_1,\dots,\mu_s$. Therefore each $c_t e_{i_tj_t}$ is a linear combination of $z_1,\dots,z_s$, so $e_{i_tj_t}\in I$ for every $t$. In particular, $I$ contains some $e_{kl}$ with $k\neq l$.

    Fix such an element $e_{kl}\in I$. Since $I$ is an ideal,
    \[
    h_{kl}:=[e_{kl},e_{lk}]=e_{kk}-e_{ll}\in I.
    \]
    For every $j\neq k$,
    \[
    [h_{kl},e_{kj}]=
    \begin{cases}
        2e_{kl}, & j=l,\\
        e_{kj}, & j\neq l,
    \end{cases}
    \]
    so $e_{kj}\in I$ for all $j\neq k$. Similarly, for every $i\neq k$,
    \[
    [e_{ik},h_{kl}]=
    \begin{cases}
        2e_{lk}, & i=l,\\
        e_{ik}, & i\neq l,
    \end{cases}
    \]
    hence $e_{ik}\in I$ for all $i\neq k$.

    Now let $i\neq j$. If $i=k$ or $j=k$, then we already know $e_{ij}\in I$. If $i,j\neq k$, then
    \[
    e_{ij}=[e_{ik},e_{kj}]\in I.
    \]
    Thus every off-diagonal matrix unit $e_{ij}$ belongs to $I$.

    Finally, for $i\neq j$,
    \[
    [e_{ij},e_{ji}]=e_{ii}-e_{jj}\in I.
    \]
    These elements span $\mathfrak{h}$, hence $\mathfrak{h}\subseteq I$. Therefore
    \[
    \mathfrak{sl}_n(F)=\mathfrak{h}\oplus \bigoplus_{i\neq j}Fe_{ij}\subseteq I,
    \]
    so $I=\mathfrak{sl}_n(F)$. Hence $\mathfrak{sl}_n(F)$ is simple.
\end{proof}

\textbf{The simplicity of $\mathfrak{o}_n$}
9`='
\begin{proof}
Let $\mathfrak{g}=\mathfrak{o}_n(F)$. For $i\neq j$, set
\[
X_{ij}=E_{ij}-E_{ji},
\]
so $X_{ji}=-X_{ij}$ and $\{X_{ij}\mid 1\le i<j\le n\}$ is a basis of $\mathfrak{g}$. The bracket is
\[
[X_{ij},X_{kl}]=\delta_{jk}X_{il}-\delta_{ik}X_{jl}-\delta_{jl}X_{ik}+\delta_{il}X_{jk}.
\]

If $n=1$ or $n=2$, then $\mathfrak{g}$ is abelian, hence not simple. For $n=4$, let
\[
I_+=\operatorname{span}\{X_{12}+X_{34},\,X_{13}-X_{24},\,X_{14}+X_{23}\},
\]
\[
I_-=\operatorname{span}\{X_{12}-X_{34},\,X_{13}+X_{24},\,X_{14}-X_{23}\}.
\]
A direct computation shows that $I_+$ and $I_-$ are ideals, $[I_+,I_-]=0$, and
\[
\mathfrak{o}_4=I_+\oplus I_-.
\]
Therefore $\mathfrak{o}_4$ is not simple.

Now assume $n=3$ or $n\ge5$, and let $I\neq0$ be an ideal of $\mathfrak{g}$. It is enough to show that $I$ contains some basis element $X_{pq}$. Indeed, if $X_{pq}\in I$ and $r\notin\{p,q\}$, then
\[
[X_{pq},X_{qr}]=X_{pr},\qquad [X_{rp},X_{pq}]=X_{rq},
\]
so repeated bracketing gives every basis element of $\mathfrak{g}$.

For $i\neq j$, define
\[
P_{ij}(x)=-[X_{ij},[X_{ij},x]].
\]
Since $I$ is an ideal, $P_{ij}(I)\subseteq I$. Moreover,
\[
P_{ij}(X_{kl})=-[X_{ij},[X_{ij},X_{kl}]]
\]
and this can be computed directly from the bracket formula.

If $|\{i,j\}\cap\{k,l\}|=0$, then all Kronecker deltas vanish, so
\[
[X_{ij},X_{kl}]=0,
\qquad\text{hence}\qquad
P_{ij}(X_{kl})=0.
\]

If $|\{i,j\}\cap\{k,l\}|=2$, then $\{k,l\}=\{i,j\}$, so $X_{kl}=\pm X_{ij}$. Hence
\[
[X_{ij},X_{kl}]=0,
\qquad\text{and again}\qquad
P_{ij}(X_{kl})=0.
\]

Now assume $|\{i,j\}\cap\{k,l\}|=1$. There are four possibilities:
\[
[X_{ij},X_{jl}]=X_{il}, \qquad
P_{ij}(X_{jl})=-[X_{ij},X_{il}]=X_{jl},
\]
\[
[X_{ij},X_{il}]=-X_{jl}, \qquad
P_{ij}(X_{il})=-[X_{ij},-X_{jl}]=X_{il},
\]
\[
[X_{ij},X_{ki}]=X_{jk}, \qquad
P_{ij}(X_{ki})=-[X_{ij},X_{jk}]=-X_{ik}=X_{ki},
\]
\[
[X_{ij},X_{kj}]=-X_{ki}, \qquad
P_{ij}(X_{kj})=-[X_{ij},-X_{ki}]=X_{kj}.
\]
Therefore
\[
P_{ij}(X_{kl})=
\begin{cases}
X_{kl}, & |\{i,j\}\cap\{k,l\}|=1, \\
0,      & |\{i,j\}\cap\{k,l\}|=0 \text{ or } 2.
\end{cases}
\]

If $n=3$, write
\[
x=c_{12}X_{12}+c_{13}X_{13}+c_{23}X_{23}\in I,\qquad x\neq0.
\]
Then
\[
x-P_{12}(x)=c_{12}X_{12},\qquad
x-P_{13}(x)=c_{13}X_{13},\qquad
x-P_{23}(x)=c_{23}X_{23}.
\]
Since not all $c_{ij}$ vanish, one of these is a non-zero scalar multiple of a basis element. Hence $I$ contains a basis element, so $I=\mathfrak{o}_3$.

Assume now $n\ge5$. Choose
\[
0\neq x=\sum_{1\le a<b\le n} c_{ab}X_{ab}\in I,
\]
and relabel indices so that $c_{12}\neq0$. Set
\[
x_1=P_{13}(x),\qquad x_2=P_{24}(x_1),\qquad x_3=P_{15}(x_2),\qquad x_4=P_{25}(x_3).
\]
By the description of $P_{ij}$,
\[
x_2\in \operatorname{span}\{X_{12},X_{14},X_{23},X_{34}\},
\]
\[
x_3\in \operatorname{span}\{X_{12},X_{14}\},
\]
and therefore
\[
x_4=c_{12}X_{12}\in I.
\]
Thus $X_{12}\in I$, hence $I=\mathfrak{o}_n$.

Therefore $\mathfrak{o}_n$ is simple exactly when $n=3$ or $n\ge5$.
\end{proof}

\textbf{The simplicity of $\mathfrak{sp}_{2n}$}

\begin{proof}
    If $n=1$, then $\mathfrak{sp}_2(F)\cong \mathfrak{sl}_2(F)$, which is simple. So assume $n\ge2$, and let
    \[
        \mathfrak{g}=\mathfrak{sp}_{2n}(F)=\left\{\begin{pmatrix} A & B \\ C & -A^T \end{pmatrix} \;\middle|\; B^T=B,\ C^T=C \right\}.
    \]
    For $1\le i,j\le n$, define
    \[
        A_{ij}=E_{ij}-E_{n+j,n+i},\qquad
        B_{ij}=E_{i,n+j}+E_{j,n+i},\qquad
        C_{ij}=E_{n+i,j}+E_{n+j,i}.
    \]
    Then $B_{ij}=B_{ji}$ and $C_{ij}=C_{ji}$, and a basis of $\mathfrak{g}$ is
    \[
        \{A_{ij}\mid i\neq j\}\cup \{H_i:=A_{ii}\mid 1\le i\le n\}\cup \{B_{ij},C_{ij}\mid 1\le i\le j\le n\}.
    \]
    Let $\mathfrak{h}=\operatorname{span}\{H_1,\dots,H_n\}$. A direct computation gives
    \[
        [H_r,A_{ij}]=(\delta_{ri}-\delta_{rj})A_{ij},
    \]
    \[
        [H_r,B_{ij}]=(\delta_{ri}+\delta_{rj})B_{ij},
        \qquad
        [H_r,C_{ij}]=-(\delta_{ri}+\delta_{rj})C_{ij}.
    \]
    Thus every basis element outside $\mathfrak{h}$ is a common eigenvector for $\operatorname{ad}\mathfrak{h}$.

    Let $I\neq0$ be an ideal of $\mathfrak{g}$. Choose $0\neq x\in I$. If $x\notin \mathfrak{h}$, then exactly as in the proof for $\mathfrak{sl}_n$, a Vandermonde argument with a suitable $h\in\mathfrak{h}$ isolates a non-zero basis vector in $I$. If $x\in \mathfrak{h}$, write
    \[
        x=\lambda_1H_1+\cdots+\lambda_nH_n,
    \]
    and choose $i$ with $\lambda_i\neq0$. Then
    \[
        [x,B_{ii}]=2\lambda_i B_{ii}\in I.
    \]
    Hence in all cases $I$ contains one of the basis elements $A_{ij}$ $(i\neq j)$, $B_{ij}$, or $C_{ij}$.

    We now show that this already implies $I=\mathfrak{g}$.

    First suppose that $B_{ii}\in I$ for some $i$. Then for every $j\neq i$,
    \[
        A_{ij}=\frac{1}{2}[B_{ii},C_{ij}]\in I,\qquad
        B_{ij}=\frac{1}{2}[A_{ij},B_{jj}]\in I,\qquad
        C_{ij}=-\frac{1}{2}[A_{ij},C_{ii}]\in I.
    \]
    Also,
    \[
        H_i=\frac{1}{4}[B_{ii},C_{ii}]\in I.
    \]
    For $p,q\neq i$ with $p\neq q$, we then have
    \[
        A_{pq}=[B_{pi},C_{iq}]\in I,\qquad
        B_{pq}=[A_{pi},B_{iq}]\in I,\qquad
        C_{pq}=-[A_{ip},C_{iq}]\in I.
    \]
    For $p\neq i$, the same computations give
    \[
        B_{pp}=[A_{pi},B_{ip}]\in I,\qquad C_{pp}=-[A_{ip},C_{ip}]\in I.
    \]
    Finally, for $p\neq i$,
    \[
        H_i+H_p=[B_{ip},C_{ip}]\in I,
    \]
    so $H_p\in I$. Therefore every basis element of $\mathfrak{g}$ lies in $I$, and hence $I=\mathfrak{g}$.

    If $B_{ij}\in I$ with $i\neq j$, then
    \[
        H_i+H_j=[B_{ij},C_{ij}]\in I,
    \]
    and hence
    \[
        B_{ii}=\frac{1}{2}[H_i+H_j,B_{ii}]\in I.
    \]
    So we again reduce to the first case.

    If $A_{ij}\in I$ with $i\neq j$, then
    \[
        B_{ij}=\frac{1}{2}[A_{ij},B_{jj}]\in I,
    \]
    and so
    \[
        H_i+H_j=[B_{ij},C_{ij}]\in I.
    \]
    Consequently,
    \[
        B_{ii}=\frac{1}{2}[H_i+H_j,B_{ii}]\in I,
    \]
    and we reduce to the previous case.

    If $C_{ii}\in I$, then
    \[
        H_i=-\frac{1}{4}[C_{ii},B_{ii}]\in I,
    \]
    so
    \[
        B_{ii}=\frac{1}{2}[H_i,B_{ii}]\in I.
    \]
    Again we reduce to the first case.

    If $C_{ij}\in I$ with $i\neq j$, then
    \[
        H_i+H_j=[B_{ij},C_{ij}]\in I,
    \]
    hence
    \[
        C_{ii}=-\frac{1}{2}[H_i+H_j,C_{ii}]\in I.
    \]
    Also,
    \[
        H_i=-\frac{1}{4}[C_{ii},B_{ii}]\in I,
    \]
    so
    \[
        B_{ii}=\frac{1}{2}[H_i,B_{ii}]\in I.
    \]
    Again we reduce to the first case.

    Therefore every non-zero ideal of $\mathfrak{g}$ is equal to $\mathfrak{g}$, so $\mathfrak{sp}_{2n}(F)$ is simple.
\end{proof}

\textbf{The simplicity of $W(n)$}

\begin{proof}
    Let
    \[
        W(n)=\operatorname{Der}\bigl(F[x_1,\dots,x_n]\bigr)
        =\left\{\sum_{i=1}^n f_i\partial_i \;\middle|\; f_i\in F[x_1,\dots,x_n]\right\},
        \qquad \partial_i=\frac{\partial}{\partial x_i},
    \]
    where $\operatorname{char}F=0$. This Lie algebra is non-abelian, since
    \[
        [\partial_i,x_i]=1\neq0.
    \]
    Let $I\neq0$ be an ideal of $W(n)$. We prove that $I=W(n)$.

    Choose
    \[
        0\neq D=\sum_{i=1}^n f_i\partial_i\in I.
    \]
    By replacing $D$ with repeated brackets $[\partial_k,D]$, we may successively differentiate the coefficients. Since $\operatorname{char}F=0$, differentiating a non-constant polynomial lowers its degree without annihilating its highest term. Hence, after finitely many such brackets, we obtain a non-zero constant vector field
    \[
        C=\sum_{i=1}^n c_i\partial_i\in I.
    \]
    Choose $k$ with $c_k\neq0$. Then for every $j$,
    \[
        [x_k\partial_j,C]
        =\sum_{i=1}^n c_i[x_k\partial_j,\partial_i]
        =-c_k\partial_j\in I.
    \]
    Therefore $\partial_j\in I$ for all $j=1,\dots,n$.

    Now let $\alpha=(\alpha_1,\dots,\alpha_n)\in \mathbb{N}^n$ and fix $j$. Since $\operatorname{char}F=0$, the scalar $\alpha_1+1$ is invertible in $F$. Set
    \[
        E=-\frac{1}{\alpha_1+1}x_1^{\alpha_1+1}x_2^{\alpha_2}\cdots x_n^{\alpha_n}\partial_j\in W(n).
    \]
    Because $\partial_1\in I$ and $I$ is an ideal,
    \[
        [E,\partial_1]
        =-\partial_1\!\left(-\frac{1}{\alpha_1+1}x_1^{\alpha_1+1}x_2^{\alpha_2}\cdots x_n^{\alpha_n}\right)\partial_j
        =x_1^{\alpha_1}\cdots x_n^{\alpha_n}\partial_j
        \in I.
    \]
    Thus every monomial derivation $x^\alpha\partial_j$ lies in $I$. These elements form a basis of $W(n)$, so $I=W(n)$.

    Therefore $W(n)$ is simple.
\end{proof}

\section{Homework week 2}

\begin{exercise}
    Prove that if $A$ is a finite-dimensional associative algebra and for every element $a \in A$ one has $a^n = 0$ for some $n$, then the algebra $A$ is nilpotent. Do not use Wedderburn's theorems. Use Engel's theorem.
\end{exercise}

\begin{proof}
    Let $A$ be a finite-dimensional associative algebra over a field, such that every element $a \in A$ is nilpotent. We want to show that $A$ is nilpotent as an algebra, meaning there exists some integer $N > 0$ such that $A^N = 0$.
    
    Consider the left regular representation of $A$ on itself. For each $a \in A$, define the left multiplication map $L_a : A \to A$ by $L_a(x) = ax$ for all $x \in A$. Since $A$ is an associative algebra, the mapping $a \mapsto L_a$ is an algebra homomorphism from $A$ to the algebra of endomorphisms $\mathrm{End}(A)$.
    
    We can view $\mathrm{End}(A)$ as a Lie algebra $\mathfrak{gl}(A)$ with the standard commutator bracket $[T, S] = T \circ S - S \circ T$. The set of all left multiplication maps $\mathfrak{g} = \{ L_a \mid a \in A \}$ forms a Lie subalgebra of $\mathfrak{gl}(A)$, because
    \[ [L_a, L_b] = L_a \circ L_b - L_b \circ L_a = L_{ab} - L_{ba} = L_{[a,b]}, \]
    and $a, b \in A \implies [a, b] \in A$, so $L_{[a,b]} \in \mathfrak{g}$.
    
    By hypothesis, for every $a \in A$, there exists an integer $n > 0$ such that $a^n = 0$. Consequently, the endomorphism $L_a$ is strictly nilpotent, since $L_a^n(x) = a^n x = 0$ for all $x \in A$. Therefore, $\mathfrak{g}$ is a Lie subalgebra of $\mathfrak{gl}(A)$ consisting entirely of nilpotent endomorphisms.
    
    According to Engel's Theorem, if a Lie subalgebra of $\mathfrak{gl}(V)$ (where $V$ is a non-zero finite-dimensional vector space) consists of nilpotent elements, then there exists a flag of subspaces
    \[ 0 = V_0 \subset V_1 \subset V_2 \subset \dots \subset V_k = A \]
    such that $\mathfrak{g}(V_i) \subseteq V_{i-1}$ for all $1 \le i \le k$.
    
    Translating this back to the associative algebra $A$, we have $L_a(V_i) \subseteq V_{i-1}$ for all $a \in A$, which is precisely equivalent to the statement $A \cdot V_i \subseteq V_{i-1}$.
    
    Using this property repeatedly, we obtain:
    \[ A \cdot V_k \subseteq V_{k-1} \]
    \[ A^2 \cdot V_k = A \cdot (A \cdot V_k) \subseteq A \cdot V_{k-1} \subseteq V_{k-2} \]
    \[ \dots \]
    \[ A^k \cdot V_k \subseteq V_0 = 0. \]
    
    Since $V_k = A$, we conclude that $A^{k+1} = A^k \cdot A \subseteq 0$. Thus, $A$ is a nilpotent associative algebra.
    \end{proof}

\begin{exercise}
    Let $\operatorname{char} F = 0$; let $L$ be a finite-dimensional solvable Lie algebra; let $e_1, \dots, e_m$ be a basis of the algebra $L$. Let $M$ be a finite-dimensional $L$-module. Suppose that every element $e_i$, for $1 \le i \le m$, acts on $M$ nilpotently. Prove that the algebra $L$ acts on $M$ nilpotently.
\end{exercise}

\begin{proof}
    Let $\bar{F}$ be the algebraic closure of $F$. We consider the scalar extensions of the Lie algebra and the module: $\bar{L} = L \otimes_F \bar{F}$ and $\bar{M} = M \otimes_F \bar{F}$.
    
    Since $L$ is a finite-dimensional solvable Lie algebra over $F$, its scalar extension $\bar{L}$ remains a solvable Lie algebra over $\bar{F}$. Furthermore, the elements $e_1, \dots, e_m$ naturally form a basis for $\bar{L}$ as a vector space over $\bar{F}$.
    
    Because $\operatorname{char} \bar{F} = 0$ and $\bar{F}$ is algebraically closed, we can apply Lie's Theorem to the finite-dimensional $\bar{L}$-module $\bar{M}$. Lie's Theorem guarantees the existence of a basis for $\bar{M}$ such that the representation $\rho: \bar{L} \to \operatorname{gl}(\bar{M})$ maps every element of $\bar{L}$ to an upper triangular matrix.
    
    Consequently, for each $x \in \bar{L}$, the diagonal entries of the matrix $\rho(x)$ are given by a set of linear functionals $\lambda_1(x), \dots, \lambda_n(x)$, where $n = \dim_{\bar{F}} \bar{M}$ and $\lambda_k \in \bar{L}^*$.
    
    We are given that each basis element $e_i$ ($1 \le i \le m$) acts nilpotently on $M$. This property is preserved under scalar extension, so each $e_i$ also acts nilpotently on $\bar{M}$. An upper triangular matrix that is nilpotent must have all its diagonal entries equal to zero. Thus, for each $i \in \{1, \dots, m\}$ and each $k \in \{1, \dots, n\}$, we have:
    \[
    \lambda_k(e_i) = 0.
    \]
    
    Since the linear functionals $\lambda_k$ vanish on the basis $e_1, \dots, e_m$, and linear transformations are determined by their action on a basis, they must vanish identically on the entire Lie algebra $\bar{L}$. That is, $\lambda_k(x) = 0$ for all $x \in \bar{L}$.
    
    This implies that for any $x \in \bar{L}$, the matrix $\rho(x)$ is strictly upper triangular. Consequently, every $x \in \bar{L}$ acts as a nilpotent endomorphism on $\bar{M}$. 
    
    Restricting this result back to the original base field $F$ and the original module $M$, we conclude that every element $x \in L$ acts nilpotently on $M$. By Engel's Theorem, a Lie algebra of nilpotent endomorphisms acts nilpotently, which completes the proof.
    \end{proof}

\begin{exercise}
    For every prime $p$, find an example showing that Lie's theorem fails.
\end{exercise}

\begin{proof}
    Let $F$ be an algebraically closed field of characteristic $p > 0$. We will construct a finite-dimensional solvable Lie algebra $\mathfrak{g}$ and a finite-dimensional representation $V$ over $F$ such that $\mathfrak{g}$ has no common eigenvector in $V$, thus showing that Lie's theorem fails.
    
    Let $V$ be a $p$-dimensional vector space over $F$ with a chosen basis $v_0, v_1, \dots, v_{p-1}$. 
    We define two linear transformations $x, y \in \mathfrak{gl}(V)$ by defining their actions on the basis vectors:
    \begin{align*}
    x(v_i) &= v_{i-1} \quad \text{for } 0 \le i \le p-1, \\
    y(v_i) &= i v_i \quad \text{for } 0 \le i \le p-1,
    \end{align*}
    where the subscripts for the action of $x$ are taken modulo $p$ (so $x(v_0) = v_{p-1}$).
    
    Let $\mathfrak{g}$ be the Lie subalgebra of $\mathfrak{gl}(V)$ generated by $x$ and $y$. First, we compute the commutator $[x, y]$ acting on an arbitrary basis vector $v_i$:
    \[
    [x, y](v_i) = x(y(v_i)) - y(x(v_i)) = x(i v_i) - y(v_{i-1}) = i v_{i-1} - (i-1) v_{i-1} = v_{i-1} = x(v_i).
    \]
    Since this holds for all basis vectors, we have $[x, y] = x$. 
    This shows that the vector space spanned by $x$ and $y$ is closed under the Lie bracket, so $\mathfrak{g} = \text{span}_F\{x, y\}$ is a 2-dimensional Lie algebra. 
    Furthermore, its derived algebra is $\mathfrak{g}' = [\mathfrak{g}, \mathfrak{g}] = \text{span}_F\{x\}$. Since $\mathfrak{g}'$ is 1-dimensional, it is abelian, which implies that $\mathfrak{g}$ is a solvable Lie algebra.
    
    Next, we show that $\mathfrak{g}$ has no common eigenvector in $V$. 
    Suppose, for the sake of contradiction, that $v = \sum_{i=0}^{p-1} c_i v_i$ is a common non-zero eigenvector for both $x$ and $y$.
    Since $v$ is an eigenvector of $y$, and the eigenvalues of $y$ are exactly $0, 1, \dots, p-1$ (which are all distinct in $F$ because the characteristic is $p$), $v$ must belong to one of the 1-dimensional eigenspaces of $y$. 
    Thus, $v$ must be a non-zero scalar multiple of $v_k$ for some $0 \le k \le p-1$.
    
    However, $v$ must also be an eigenvector for $x$. We check the action of $x$ on $v_k$:
    \[
    x(v_k) = v_{k-1} \pmod p.
    \]
    Because $p$ is a prime (and thus $p \ge 2$), $v_{k-1}$ is linearly independent from $v_k$. Therefore, $v_k$ is not an eigenvector for $x$, which is a contradiction.
    
    Consequently, the solvable Lie algebra $\mathfrak{g}$ has no common eigenvector in the representation $V$, demonstrating that Lie's theorem does not hold in characteristic $p$.
    \end{proof}




















\end{document}
